% 中文译稿；底本 /tmp/sga3-audit/units/XV.tex，源行 1--173。
\markboth{第 XV 讲. 子环面的补充}{第 XV 讲. 子环面的补充}
\chapter*{第 XV 讲\\群预概形的子环面之补充\\在光滑群中的应用}
\addcontentsline{toc}{chapter}{第 XV 讲. 群预概形的子环面之补充}
\SGARefTarget{sga3:xv:expose}

\begin{center}
\textit{M. Raynaud}\footnote{底本版本注：xy 版本日期为 2008 年 12 月 1 日。}\textsuperscript{,}\footnote{本讲及随后两讲（第 XVI、XVII 讲）并非讨论班口头报告的讲稿。它们以 A. Grothendieck 于 1964 年为本讨论班撰写的简略手稿笔记为基础，加以展开并补充了一些材料。}
\end{center}

\SGASection{0}{引言}
\SGARefTarget{sga3:xv:section:0}

本讲补充并部分改写第 XI、XII 讲；不必预先掌握第 XIII、XIV 讲。我们继续第 XII 讲开始的研究，考虑不一定仿射、也不一定在 \(S\) 上分离的 \(S\)-群预概形。

\SGARefLink{sga3:xv:section:1}{第 1 节}以及第 2、3、4 节研究群预概形的子环面。我们将得到无穷小提升定理（第 2 节）和整体提升定理（第 4 节），其中有限阶点（第 1 节）起关键作用。

\SGARefLink{sga3:xv:section:5}{第 5 节}以及第 6、7 节与前面各节相互独立。通过考虑无穷小邻域，可证明那些等于其连通正规化子的光滑子群所成函子可表（第 5 节）。第 6、7 节进一步专门研究 Cartan 子群。

最后，第 8 节给出光滑群的子环面函子或极大环面函子可表的一个充要条件。

\SGASection{1}{有限子群的提升}
\SGARefTarget{sga3:xv:section:1}

\SGASubsection{1.1}{有限、光滑、中心的乘法型子群}
\SGARefTarget{sga3:xv:subsection:1:1}

\SGAStatement{命题 1.1}{设 \(S\) 为仿射概形，\(S_0\) 为 \(S\) 中由平方为零的理想定义的闭子概形，\(G\) 为 \(S\)-群预概形，\(H_0\) 为 \(G_0=G\times_S S_0\) 的子群概形，且在 \(S_0\) 上光滑、有限型并为乘法型。存在提升 \(H_0\) 的 \(G\) 的乘法型子群概形，当且仅当存在提升 \(H_0\) 的 \(G\) 的子概形 \(H'\)，使其在 \(S\) 上平坦。}
\SGARefTarget{sga3:xv:proposition:1:1}

条件的必要性显然；下面证明充分性。乘法型群 \(H_0\) 是拟等平凡的（\SGARefLink{sga3:x:corollary:4-5}{X~4.5}）。由第 \SGARefLink{sga3:x:proposition:2:1}{X~2.1} 讲，存在乘法型 \(S\)-群 \(H\) 以及群的一个 \(S_0\)-同构
\[
  u_0:H\times_S S_0\xrightarrow{\sim}H_0.
\]
由于 \(H'\) 在 \(S\) 上平坦，且 \(H'\times_S S_0\) 在 \(S_0\) 上有限表现（第 \SGARefLink{sga3:ix:proposition:2:1:item:b:01}{IX~2.1(b)} 讲），\(H'\) 在 \(S\) 上有限表现。其纤维光滑，故 \(H'\) 在 \(S\) 上光滑（EGA IV~17.5.1）。由于 \(S\) 仿射，\(u_0\) 因而可提升为预概形的一个 \(S\)-态射
\[
  u:H\longrightarrow H'.
\]
由第 \SGARefLink{sga3:iii:theorem:2-1}{III~2.1} 讲和第 \SGARefLink{sga3:ix:theorem:3:1}{IX~3.1} 讲，复合态射
\[
  v_0:H\times_S S_0\xrightarrow{\sim}H_0\longrightarrow G_0
\]
也可提升为一个 \(S\)-同态
\[
  v:H\longrightarrow G.
\]
由于 \(v_0\) 是浸入，\(v\) 也是浸入。因此 \(H\) 在 \(v\) 下的像是 \(G\) 的一个乘法型子群概形，并提升了 \(H_0\)。

\SGAStatement{命题 1.2}{设 \(S\) 为预概形，\(S_0\) 为 \(S\) 中由局部幂零理想层定义的子预概形，\(G\) 为在 \(S\) 上平坦且有限表现的 \(S\)-群预概形，\(H_0\) 为 \(G_0=G\times_S S_0\) 的子群概形，且在 \(S_0\) 上光滑、有限、为乘法型并包含在中心中。则存在唯一的提升 \(H_0\) 的 \(G\) 的乘法型子群概形 \(H\)。此外，\(H\) 包含在中心中。（见 XVII，附录 III，1。）}
\SGARefTarget{sga3:xv:proposition:1:2}

\SGAInlineRefTarget{sga3:xv:proposition:1-2-bis}\SGAStatement{命题 1.2 bis}{设 \(S,G,S_0,G_0\) 如上，\(H\) 为在 \(S\) 上光滑且有限的乘法型 \(S\)-群概形，\(u_0:H_0=H\times_S S_0\to G_0\) 为中心同态。则 \(u_0\) 唯一提升为同态 \(u:H\to G\)。此外，\(u\) 是中心同态。}
\SGARefTarget{sga3:xv:proposition:1:2:bis}

考虑 \(u_0\) 的图，由\SGARefLink{sga3:xv:proposition:1:2}{命题 1.2} 即得\SGARefLink{sga3:xv:proposition:1:2:bis}{命题 1.2 bis} 中提升 \(u\) 的存在性。由第 \SGARefLink{sga3:ix:corollary:3:4}{IX~3.4} 讲和第 \SGARefLink{sga3:ix:theorem:5:1}{IX~5.1} 讲，提升 \(u\) 唯一且为中心同态。

\SGAProof{\SGARefLink{sga3:xv:proposition:1:2}{命题 1.2} 的证明}
\(H\) 的唯一性以及 \(H\) 包含在中心中的事实，均由第 \SGARefLink{sga3:ix:theorem:5:6:item:b}{IX~5.6(b)} 讲和第 \SGARefLink{sga3:ix:corollary:3:4:bis}{IX~3.4 bis} 讲得出。由唯一性，为证明存在性，可假设 \(S\) 仿射且 \(S_0\) 由平方为零的理想定义；根据\SGARefLink{sga3:xv:proposition:1:1}{命题 1.1}，只须找到 \(G\) 中一个在 \(S\) 上平坦且提升 \(H_0\) 的子概形。

由于 \(H_0\) 在 \(S_0\) 上光滑且有限，必要时缩小 \(S\)，可假设存在在 \(S\) 上可逆的整数 \(n>0\)，使 \(H_0={}_nH_0\)。考虑 \(G\) 上的 \(n\) 次幂态射：
\[
  u:G\longrightarrow G,\qquad x\longmapsto x^n.
\]
仍用 \({}_nG\) 表示“\(u\) 的核”，即 \(G\) 的单位截面在 \(u\) 下的逆像。注意，\(u\) 一般不是群同态。我们先承认如下引理。

\SGAStatement{引理 1.3}{设 \(k\) 为域，\(G\) 为在 \(k\) 上局部有限型的群概形，\(n\) 为与 \(k\) 的特征互素的整数，\(u\) 为 \(G\) 上的 \(n\) 次幂态射。则 \(u\) 在 \(G\) 中属于 \(G\) 的中心的每一点 \(x\) 处都是平展的。}
\SGARefTarget{sga3:xv:lemma:1:3}

由于 \(G\) 在 \(S\) 上平坦且有限表现，上述引理与 EGA IV~17.8.2 表明：若 \(x\in G\) 位于 \(s\in S\) 之上且属于 \(G_s\) 的中心，则 \(u\) 在 \(x\) 处平展。若再有 \(x\in {}_nG\)，则 \({}_nG\) 在 \(x\) 处于 \(S\) 上平展。根据假设，\(H_0\) 包含在中心中且包含于 \({}_nG_0\)；因而包含于 \({}_nG\) 中在 \(S\) 上平展的最大开子集 \(V\)。由于 \(H_0\) 与 \(V\times_S S_0=V_0\) 均在 \(S_0\) 上平展，\(H_0\) 是 \(V_0\) 的开子概形（EGA IV~17.8.7、17.9.1）。因此，\(V\) 中与 \(H_0\) 有相同底空间的开子预概形，就是 \(G\) 中在 \(S\) 上平坦且提升 \(H_0\) 的子概形。

尚须证明\SGARefLink{sga3:xv:lemma:1:3}{引理 1.3}。\(G\) 中 \(u\) 平展的最大开子集在基域扩张下不变（EGA IV~17.8.2），故可归结到 \(x\) 为 \(k\)-有理点的情形。用 \(t_x\) 表示由 \(x\) 定义的平移，它是 \(G\) 的一个 \(k\)-自同构。由于 \(x\) 属于 \(G\) 的中心，图
\[
\begin{tikzcd}[column sep=large,row sep=large]
G \arrow[r,"u"] \arrow[d,"t_x"'] & G \arrow[d,"t_x^n"] \\
G \arrow[r,"u"] & G .
\end{tikzcd}
\]
交换。因此只须证明 \(u\) 在单位元处平展；这已在 VII\textsubscript{A}~\SGARefLink{sga3:xv:definition:8-4}{8.4} 中证明。

% 底本源行 174--502。
\SGASubsection{1.2}{有限群的整体提升}
\SGARefTarget{sga3:xv:subsection:1:2}

\SGAStatement{\SGAInlineRefTarget{sga3:xv:lemma:1-4}引理 1.4}{设 \(A\) 为局部环，关于其极大理想 \(\mathfrak m\) 定义的拓扑分离且完备，置
\[
  S=\operatorname{Spec}(A),\qquad
  S_n=\operatorname{Spec}(A/\mathfrak m^n).
\]
则对每个预概形 \(X\)（分别地，对每个 \(S\)-预概形 \(X\)），典范映射
\[
  \tag{$\ast$}\label{sga3:xv:lemma:1-4:equation:ast}
  \operatorname{Hom}(S,X)
  \longrightarrow
  \varprojlim_n\operatorname{Hom}(S_n,X)
\]
（分别地，
\[
  \tag{$\ast\ast$}\label{sga3:xv:lemma:1-4:equation:ast-ast}
  \Gamma(X/S)
  \longrightarrow
  \varprojlim_n\Gamma(X_n/S_n),
  \qquad X_n=X\times_S S_n)
\]
为双射。}
\SGARefTarget{sga3:xv:lemma:1:4}

断言 \((\ast\ast)\) 容易由 \((\ast)\) 推出。下面证明 \((\ast)\)。设
\[
  u_n:S_n\longrightarrow X\qquad(n\in\mathbb N)
\]
为一族相容的态射。\(S_n\) 的闭点的像 \(y\) 与 \(n\) 无关，且 \(u_n\) 通过 \(\operatorname{Spec}\mathcal O_y\) 分解。取逆极限，态射 \(u_n\) 定义环同态
\[
  \widetilde u:\mathcal O_y
  \longrightarrow
  \varprojlim_n(A/\mathfrak m^n)=A.
\]
这证明 \((\ast)\) 是满射；只要 \(A\) 关于 \(\mathfrak m\)-进拓扑分离，此映射就是单射。

\SGAStatement{\SGAInlineRefTarget{sga3:xv:corollary:1-5}推论 1.5}{设 \(A\) 为完备 noether 局部环，\(\mathfrak m\) 为其极大理想，置
\[
  S=\operatorname{Spec}(A),\qquad
  S_n=\operatorname{Spec}(A/\mathfrak m^n),
\]
设 \(X\) 为 \(S\) 上的有限概形，\(Y\) 为 \(S\)-预概形。则典范映射
\[
  \operatorname{Hom}_S(X,Y)
  \longrightarrow
  \varprojlim_n\operatorname{Hom}_{S_n}(X_n,Y_n)
\]
为双射（其中 \(X_n=X\times_S S_n\)，\(Y_n\) 同理）。}
\SGARefTarget{sga3:xv:corollary:1:5}

事实上，由 EGA II~6.2.5，\(X\) 是有限个在 \(S\) 上有限的局部 \(S\)-概形的不交和。因而只须处理 \(X\) 本身是一个完备 noether 局部环的谱的情形。但
\[
  \operatorname{Hom}_S(X,Y)=\Gamma(Z/X),
  \qquad Z=Y\times_S X,
\]
故可应用上述引理。

\SGAStatement{\SGAInlineRefTarget{sga3:xv:proposition:1-6}命题 1.6}{设 \(A,S,S_n\) 如上，\(G\)、\(M\) 为两个 \(S\)-群预概形，且 \(M\) 在 \(S\) 上有限。则：
\begin{enumerate}
\item[\normalfont a)]\SGAInlineRefTarget{sga3:xv:proposition:1-6:item:a} 典范映射
\[
  \operatorname{Hom}_{S\text{-}\mathrm{gr}}(M,G)
  \longrightarrow
  \varprojlim_n
  \operatorname{Hom}_{S_n\text{-}\mathrm{gr}}(M_n,G_n)
\]
为双射。
\item[\normalfont b)] 若 \(M\) 为乘法型且 \(G\) 在 \(S\) 上光滑，则典范映射
\[
  \varphi:\operatorname{Hom}_{S\text{-}\mathrm{gr}}(M,G)
  \longrightarrow
  \operatorname{Hom}_{S_0\text{-}\mathrm{gr}}(M_0,G_0)
\]
为满射。此外，若 \(\varphi(u)=\varphi(u')=u_0\)，则 \(u\)、\(u'\) 可由 \(G(S)\) 中一个约化为 \(G(S_0)\) 的单位元的元素共轭。
\item[\normalfont c)] 若 \(M\) 为乘法型且在 \(S\) 上光滑，\(G\) 在 \(S\) 上平坦且有限型，且 \(u_0:M_0\to G_0\) 为中心同态，则 \(u_0\) 唯一提升为同态
\[
  u:M\longrightarrow G.
\]
此外，若 \(G\) 的纤维连通，则 \(u\) 是中心同态。
\end{enumerate}}
\SGARefTarget{sga3:xv:proposition:1:6}

\SGAProof{证明}
\emph{a)} 这由 \SGARefLink{sga3:xv:corollary:1-5}{1.5}、\(M\times_S M\) 在 \(S\) 上有限，以及群态射可刻画为使下列熟知的图交换的态射而得：
\[
\begin{tikzcd}[column sep=large,row sep=large]
M\times_S M \arrow[r,"u\times u"] \arrow[d]
  & G\times_S G \arrow[d] \\
M \arrow[r,"u"'] & G
\end{tikzcd}
\]

\emph{b)} 根据第 \SGARefLink{sga3:ix:theorem:3:6}{IX~3.6} 讲，可以构造一族相容的同态 \(u_n:M_n\to G_n\)，提升给定同态 \(u_0:M_0\to G_0\)。结合 a)，这就证明了 b) 的第一个断言。

若 \(u\)、\(u'\) 为 \(u_0\) 的两个提升，则 \(u_n\)、\(u'_n\) 可由一个提升 \(G(S_0)\) 的单位元的元素 \(g_n\in G(S_n)\) 共轭（第 \SGARefLink{sga3:ix:theorem:3:6}{IX~3.6} 讲）。同一结果还表明可以相容地选取 \(g_n\)；故由 \SGARefLink{sga3:xv:lemma:1-4}{1.4}，它们来自一个截面 \(g\in G(S)\)。于是态射 \(u\)、\(\operatorname{Int}(g)u'\) 对每个 \(n\) 都模 \(\mathfrak m^n\) 相同，因而由 \SGARefLink{sga3:xv:corollary:1-5}{1.5}，两者相同。

\emph{c)} \(u\) 的存在性与唯一性由 a) 和 \SGARefLink{sga3:xv:proposition:1:2:bis}{1.2 bis} 推出。若 \(G\) 的纤维连通，则由第 \SGARefLink{sga3:ix:theorem:5:6:item:a}{IX~5.6(a)} 讲，\(u\) 是中心同态。

\SGASection{2}{子环面的无穷小提升}
\SGARefTarget{sga3:xv:section:2}

\SGASubsection{2.1}{定理的陈述}
\SGARefTarget{sga3:xv:subsection:2:1}

我们将给出群预概形中子环面的一个无穷小提升定理，它不要求光滑性假设（与第 \SGARefLink{sga3:ix:theorem:3:6:bis}{IX~3.6 bis} 讲不同），同时回答一个十分自然的问题：\footnote{编者注：用环面的有限子群逼近该环面的想法，已见于 Grothendieck 关于环面的中心化子之连通性的证明；参见 \emph{BIBLE} 第 \SGARefLink{sga3:xv:lemma:4:6}{4.6} 节。}若子环面 \(T_0\) 的“足够多”的有限阶点可以提升，是否就足以保证 \(T_0\) 本身也可以提升（当然是无穷小地提升）？

\SGAStatement{定理 2.1}{设 \(S\) 为 noether 仿射概形，\(S_0\) 为 \(S\) 中由平方为零的理想 \(J\) 定义的闭子概形，\(G\) 为有限型 \(S\)-群预概形，\(G_0=G\times_S S_0\)，\(T_0\) 为 \(G_0\) 的子环面，\(q>0\) 为在 \(S\) 上可逆的整数。假设对每个等于 \(q\) 的幂的整数 \(n\)，存在 \(G\) 的子概形 \(M_n\)，使其在 \(S\) 上平坦且满足
\[
  M_n\times_S S_0={}_nT_0.
\]
则存在 \(G\) 的子环面 \(T\)，使 \(T\times_S S_0=T_0\)。}
\SGARefTarget{sga3:xv:theorem:2:1}

\SGARefLink{sga3:xv:theorem:2:1}{定理 2.1} 将通过下述两个推论发挥作用。

\SGAStatement{\SGAInlineRefTarget{sga3:xv:corollary:2-2}推论 2.2}{设 \(S\) 为局部 noether 预概形，\(S_0\) 为 \(S\) 中由局部幂零理想层定义的闭子预概形，\(G\) 为有限型 \(S\)-群预概形，\(T_0\) 为 \(G_0=G\times_S S_0\) 的子环面，\(q>0\) 为在 \(S\) 上可逆的整数。当 \(n\) 遍取 \(q\) 的幂时，设 \((M_n)\) 为 \(G\) 的一族相容的乘法型 \(S\)-子群概形，提升各 \({}_nT_0\)。（所谓族 \((M_n)\) 相容，是指当 \(m\) 整除 \(n\) 时，
\[
  M_m={}_m(M_n)
\]
成立。）则存在唯一的 \(G\) 的子环面 \(T\)，使对每个 \(n\)，
\[
  T\times_S S_0=T_0
  \qquad\text{且}\qquad
  {}_nT=M_n
\]
成立。}
\SGARefTarget{sga3:xv:corollary:2:2}

\SGAStatement{推论 2.3}{设 \(G\) 为在 \(S\) 上平坦且有限表现的 \(S\)-群预概形，\(S_0\) 为 \(S\) 中由有限型局部幂零理想层定义的闭子预概形，\(T_0\) 为 \(G_0=G\times_S S_0\) 的中心环面。则存在唯一的 \(G\) 的子环面 \(T\) 提升 \(T_0\)。此外，\(T\) 包含在中心中。}
\SGARefTarget{sga3:xv:corollary:2:3}

\SGAStatement{注 2.4}{请读者自行表述断言 \SGARefLink{sga3:xv:theorem:2:1}{2.1}、\SGARefLink{sga3:xv:corollary:2-2}{2.2}、\SGARefLink{sga3:xv:corollary:2:3}{2.3} 的类似形式：不再提升 \(G_0\) 的子环面，而是给定 \(S\) 上的环面 \(T\)，求态射
\[
  u_0:T_0\longrightarrow G_0
\]
的提升。考虑 \(u_0\) 的图，即可归结到前述情形。}
\SGARefTarget{sga3:xv:remark:2:4}

下面说明如何由\SGARefLink{sga3:xv:theorem:2:1}{定理 2.1} 推出\SGARefLink{sga3:xv:corollary:2-2}{推论 2.2} 和 \SGARefLink{sga3:xv:corollary:2:3}{2.3}。

\SGAProof{\SGARefLink{sga3:xv:corollary:2-2}{推论 2.2} 的证明}
\(T\) 的唯一性由第 \SGARefLink{sga3:ix:corollary:4:8:item:b}{IX~4.8(b)} 讲和第 \SGARefLink{sga3:ix:remark:4:10}{IX~4.10} 讲推出。因此，为证明存在性，可归结到 \(S\) 仿射（从而为 noether）、且 \(S_0\) 由平方为零的理想定义的情形。

\SGAStatement{引理 2.5}{设 \(G\) 为有限表现 \(S\)-群预概形，\(H\) 为 \(G\) 的乘法型子群概形，且在 \(S\) 上光滑。则 \(\operatorname{Centr}_G(H)\) 可由 \(G\) 的一个有限表现 \(S\)-子群预概形表示。}
\SGARefTarget{sga3:xv:lemma:2:5}

当 \(G\) 在 \(S\) 上分离时，由第 \SGARefLink{sga3:viii:examples:6:5:item:e:01}{VIII~6.5(e)} 讲即可得到此引理，无须对 \(H\) 作任何光滑性假设。在目前的情形，注意断言关于 fpqc 拓扑是局部的。因此可以假设 \(H\) 可对角化，因而形如 \(D_S(M)\)。还可假设 \(S\) 仿射，再由 EGA IV~8 假设其为 noether。由于 \(H\) 在 \(S\) 上光滑且有限型，\(M\) 的挠子群的阶 \(q\) 在 \(S\) 上可逆（第 \SGARefLink{sga3:viii:proposition:2:1:item:e:01}{VIII~2.1(e)} 讲）。于是立即可见（参见第 \SGARefLink{sga3:ix:remark:4:10}{IX~4.10} 讲），当 \(n\) 遍取 \(q\) 的幂时，子群 \({}_nH\) 在 \(H\) 中概形稠密（第 \SGARefLink{sga3:ix:definition:4:1}{IX~4.1} 讲）。但 \({}_nH\) 是 \(S\) 的完全分裂覆盖，即同构于有限个 \(S\) 的副本的不交和。因此
\[
  \operatorname{Centr}_G({}_nH)=Z_n
\]
可表为 \({}_nH\) 在 \(S\) 上各截面在 \(G\) 中的中心化子之交。接下来只须应用如下引理。

\SGAStatement{\SGAInlineRefTarget{sga3:xv:lemma:2-5-bis}引理 2.5 bis}{设 \(S\) 为 noether 预概形，\(G\) 为有限型 \(S\)-群预概形，\(H\) 为 \(G\) 的乘法型子群，\((H_i)\) 为 \(G\) 的一族递增的乘法型子群滤过。假设
\[
  Z_i=\operatorname{Centr}_G(H_i)
\]
可由 \(G\) 的子群预概形表示。则族 \((Z_i)\) 最终稳定。若 \(H_i\) 还在 \(H\) 中概形稠密，则对充分大的 \(i\)，
\[
  Z_i=\operatorname{Centr}_G(H)
\]
成立。}
\SGARefTarget{sga3:xv:lemma:2:5:bis}

为证明族 \((Z_i)\) 最终稳定，只须证明底集族 \(\operatorname{ens}(Z_i)\) 最终稳定。事实上，稳定值此时是 \(G\) 的某开子集 \(U\) 中的闭子集。以 \(U\) 替换 \(G\) 后，即归结到研究 noether 预概形中一族递减的闭子预概形滤过。一个初等的可构造性论证可归结到 \(S\) 整的情形。此时只须证明 \(\operatorname{ens}(Z_i)\) 在 \(S\) 的某个非空开子集上最终稳定。\(G\) 的一般纤维是分离的（第 VI\textsubscript{A}~0.3 讲）；因此缩小 \(S\) 后，可假设 \(G\) 在 \(S\) 上分离（EGA IV~8）。于是 \(Z_i\) 在 \(G\) 中闭（第 VIII~6.5(e) 讲）。

为证明引理的最后一个断言，用 \(Z\) 表示族 \((Z_i)\) 的稳定值。显然，\(\operatorname{Centr}_G(H)\) 是 \(Z\) 的子函子；下面证明 \(Z\) 中的元素与 \(H\) 中的元素交换。设 \(E\) 为 \(H\times_S Z\) 的子预概形，定义为下面一对态射的核：
\[
\begin{aligned}
  H\times_S Z&\rightrightarrows G,\\
  (h,c)&\longmapsto c,\\
  (h,c)&\longmapsto hch^{-1}.
\end{aligned}
\]
对每个 \(i\)，预概形 \(E\) 都包含 \(H_i\times_S Z\)。各 \(H_i\) 在 \(S\) 上平坦，故 EGA IV~11.10.9 表明，对 \(S\) 的每一点 \(s\)，各 \((H_i)_s\) 在 \(H_s\) 中概形稠密，且各 \(H_i\times_S Z\) 在 \(H\times_S Z\) 中概形稠密。由于 \(G_s\) 分离，\(E_s\) 在 \((H\times_S Z)_s\) 中闭，因而等于后者。故 \(E\) 在 \(H\times_S Z\) 中闭，从而等于 \(H\times_S Z\)。这就是说，\(Z\) 中的元素与 \(H\) 中的元素交换。

回到 \SGARefLink{sga3:xv:corollary:2-2}{2.2} 的证明，\SGARefLink{sga3:xv:lemma:2:5}{引理 2.5} 表明
\[
  Z_n=\operatorname{Centr}_G(M_n)
\]
可表，而\SGARefLink{sga3:xv:lemma:2-5-bis}{引理 2.5 bis} 表明递减的子预概形族 \((Z_n)\) 最终稳定；用 \(Z\) 表示其稳定值。群 \(Z\) 包含 \(T_0\) 及所有 \(M_n\)。因此用 \(Z\) 替换 \(G\) 后，可假设各 \(M_n\) 包含在中心中。

此时\SGARefLink{sga3:xv:theorem:2:1}{定理 2.1} 的假设均满足，故存在 \(G\) 的子环面 \(T\) 提升 \(T_0\)。群 \({}_nT\)、\(M_n\) 都是 \({}_nT_0\) 的提升，故彼此共轭（第 \SGARefLink{sga3:ix:theorem:3:2:bis}{IX~3.2 bis} 讲）；由于 \(M_n\) 包含在中心中，两者相同。这个环面 \(T\) 即具有所需的性质。

\SGAProof{\SGARefLink{sga3:xv:corollary:2:3}{推论 2.3} 的证明}
\(T\) 的唯一性由第 \SGARefLink{sga3:ix:theorem:5:1:bis}{IX~5.1 bis} 讲推出，其中心性则由第 \SGARefLink{sga3:ix:theorem:5:6:item:b}{IX~5.6(b)} 讲推出。按通常的步骤，由此可先归结到 \(S\) 仿射（从而 \(S_0\) 由有限型幂零理想定义）的情形，再归结到 noether 情形。必要时缩小 \(S\)，可假设存在在 \(S\) 上可逆的整数 \(q\)。于是\SGARefLink{sga3:xv:corollary:2:3}{推论 2.3} 由 \SGARefLink{sga3:xv:corollary:2-2}{2.2} 和 \SGARefLink{sga3:xv:proposition:1:2}{1.2} 得到。

\SGAStatement{注 2.6}{若以 \(G_0\) 的一个光滑、中心的乘法型子群替换环面 \(T_0\)，则\SGARefLink{sga3:xv:corollary:2:3}{推论 2.3} 仍然成立。}
\SGARefTarget{sga3:xv:remark:2:6}

% 底本源行 503--848。
\SGASubsection{2.2}{\SGARefLink{sga3:xv:theorem:2:1}{定理 2.1} 的证明}
\SGARefTarget{sga3:xv:subsection:2:2}

\medskip\noindent\textbf{a) 归结到 \(\boldsymbol{T_0=\mathbb G_{m,S_0}}\) 的情形。}

由 \SGARefLink{sga3:xv:proposition:1:1}{1.1}，可假设 \(M_n\) 为乘法型子群。应用第 \SGARefLink{sga3:ix:theorem:3:2:bis}{IX~3.2 bis} 讲，可假设族 \((M_n)\) 相容（\SGARefLink{sga3:xv:corollary:2-2}{2.2}）。\(M_n\) 在 \(G\) 中的中心化子 \(Z_n\) 可表（\SGARefLink{sga3:xv:lemma:2:5}{2.5}），且滤过族 \((Z_n)\) 最终稳定（\SGARefLink{sga3:xv:lemma:2-5-bis}{2.5 bis}）。因此，取充分大的 \(n\)，用 \(Z_n\) 替换 \(G\) 后，可假设 \(T_0\) 与各 \(M_n\) 包含在中心中。此时第 \SGARefLink{sga3:ix:theorem:5:1:bis}{IX~5.1 bis} 讲保证 \(T_0\) 的提升唯一。

与 \SGARefLink{sga3:xv:proposition:1:1}{1.1} 的证明相同，可假设存在一个 \(S\)-环面 \(T\) 及一个 \(S_0\)-同构
\[
  u_0:T\times_S S_0\xrightarrow{\sim}T_0.
\]
提升 \(u_0\) 等价于提升 \(T_0\)。由唯一性，只须在一个有限型仿射忠实平坦扩张 \(S'\to S\) 之后证明 \(u_0\) 的提升存在（fpqc 下降）。于是可假设
\[
  T=\mathbb G_{m,S}^{\,r}
\]
（第 \SGARefLink{sga3:x:corollary:4-5}{X~4.5} 讲）。若 \(u_0\) 在每个因子 \(\mathbb G_m\) 上的限制均可提升为一个 \(S\)-态射——它必为中心同态——则这些提升立即给出 \(u_0\) 的一个提升。总之，可假设 \(T_0=\mathbb G_{m,S_0}\)。

\medskip\noindent\textbf{b) \(\boldsymbol{T_0}\) 的提升存在性之障碍的定义。}

由 \SGARefLink{sga3:xv:proposition:1:1}{1.1}，为证明\SGARefLink{sga3:xv:theorem:2:1}{定理 2.1}，只须找到 \(G\) 中一个在 \(S\) 上平坦且提升 \(T_0\) 的子概形。我们将看到，这种提升的存在性之障碍可定义为某个模层的 \(\operatorname{Ext}^1\) 群的元素。

取 \(G\) 中的开子集 \(U\)，使 \(T_0\) 在其中闭，仍用 \(U\)（分别地，\(U_0\)）表示 \(G\)（分别地，\(G_0\)）中以此开子集为底空间的开子概形。因此，将层 \(\mathcal O_{T_0}\) 视为 \(U\) 上的层时，它是 \(\mathcal O_{U_0}\) 的商。设
\[
  h:\mathcal O_{U_0}\longrightarrow\mathcal O_{T_0}
\]
为典范满态射。

\SGAStatement{\SGAInlineRefTarget{sga3:xv:lemma:2-7}引理 2.7}{典范映射
\[
  h'=\operatorname{id}_J\otimes h:
  J\otimes_{S_0}\mathcal O_{U_0}
  \longrightarrow
  J\otimes_{S_0}\mathcal O_{T_0}
\]
可分解为 \(i\circ j_U\)，且这种分解必然唯一，其中 \(j_U\) 为典范满态射
\[
  J\otimes_{S_0}\mathcal O_{U_0}
  \longrightarrow
  J\mathcal O_U\simeq J\mathcal O_{U_0}.
\]}
\SGARefTarget{sga3:xv:lemma:2:7}

须证明 \(h'(K)=0\)，其中 \(K\) 为 \(j_U\) 的核。对每个等于 \(q\) 的幂的整数 \(n\)，有满态射
\[
  h_n:\mathcal O_{T_0}\longrightarrow\mathcal O_{{}_nT_0}.
\]
由于 \({}_nT_0\) 可提升为 \(S\) 上平坦的概形 \(M_n\)，典范态射
\[
  j_n:J\otimes_{S_0}\mathcal O_{{}_nT_0}
  \xrightarrow{\sim}J\mathcal O_{M_n}
\]
为同构。交换图
\[
\begin{tikzcd}[column sep=large,row sep=large]
K \arrow[r]
  & J\otimes_{S_0}\mathcal O_{U_0}
      \arrow[r,"j_U"] \arrow[d,"h'"]
  & J\mathcal O_U \arrow[r,hook] \arrow[dl,"i"]
  & \mathcal O_U \\
  & J\otimes_{S_0}\mathcal O_{T_0}
      \arrow[d,"h'_n"] & & \\
  & J\otimes_{S_0}\mathcal O_{{}_nT_0}
      \arrow[r,"j_n"',"\sim"] & J\mathcal O_{M_n}
      \arrow[r,hook] & \mathcal O_{M_n}
\end{tikzcd}
\]
表明，对每个 \(n\)，\(h'(K)\) 都包含在 \(\ker h'_n\) 中。因此它包含在 \(\bigcap_n\ker h'_n\) 中，只须证明此交为零。现有
\[
  \mathcal O_{T_0}=\mathcal O_{S_0}[\mathbb Z],
  \qquad
  \mathcal O_{{}_nT_0}
     =\mathcal O_{S_0}[\mathbb Z/n\mathbb Z].
\]
设
\[
  a=\sum_{m\in\mathbb Z}a_m\otimes m
  \in J\otimes_{S_0}\mathcal O_{T_0}.
\]
其中各 \(a_m\) 为 \(J\) 的截面，除有限个外均为零。取充分大的 \(n\)，使满足 \(a_m\ne0\) 的各指标 \(m\) 在 \(\mathbb Z/n\mathbb Z\) 中的像互不相同。若 \(a\in\ker h'_n\)，则必有 \(a=0\)。故 \(\bigcap_n\ker h'_n=0\)，这就证明了 \SGARefLink{sga3:xv:lemma:2-7}{2.7}。

设 \(K_0\) 为 \(h\) 的核。考虑图
\[
\begin{tikzcd}[column sep=large,row sep=large]
  &&& 0 \arrow[d] \\
  & J\mathcal O_U \arrow[d,"\sim"] && K_0 \arrow[d] \\
0 \arrow[r]
  & J\mathcal O_{U_0} \arrow[r] \arrow[d,"i"']
  & \mathcal O_U \arrow[r] & \mathcal O_{U_0}
      \arrow[r] \arrow[d,"h"'] & 0 \\
  & J\otimes_{S_0}\mathcal O_{T_0}
  && \mathcal O_{T_0} \arrow[d] \\
  &&& 0 .
\end{tikzcd}
\]
层 \(\mathcal O_U\) 定义一个元素
\[
  \Phi\in
  \operatorname{Ext}^{1}_{\mathcal O_U}
  (\mathcal O_{U_0},J\mathcal O_{U_0}).
\]
利用 \(\operatorname{Ext}^{1}\) 的双函子性，通过态射
\[
  K_0\longrightarrow\mathcal O_{U_0},
  \qquad
  i:J\mathcal O_{U_0}
    \longrightarrow J\otimes_{S_0}\mathcal O_{T_0},
\]
由 \(\Phi\) 得到元素 \(\Psi\)，它属于
\[
  \operatorname{Ext}^{1}_{\mathcal O_U}
  (K_0,J\otimes_{S_0}\mathcal O_{T_0}).
\]
由第 \SGARefLink{sga3:iii:lemma:4-1}{III~4.1} 讲和平坦性的无穷小判据（参见第 \SGARefLink{sga3:iii:proposition:4-3}{III~4.3} 讲）可知：存在 \(U\) 中在 \(S\) 上平坦且提升 \(T_0\) 的子概形，当且仅当 \(\Psi=0\)。

由于 \(T_0\) 仿射，第 \SGARefLink{sga3:iii:proposition:4-5}{III~4.5} 讲和 \SGARefLink{sga3:xv:lemma:4:6}{4.6} 表明，只须在 \(U\) 上局部构造一个在 \(S\) 上平坦且提升 \(T_0\) 的子概形。因此只须证明 \(\Psi\) 在层
\[
  \mathscr E=
  \mathcal Ext^{1}_{\mathcal O_U}
  (K_0,J\otimes_{S_0}\mathcal O_{T_0})
\]
中的像为零。仍用 \(\Psi\) 表示 \(\Gamma(U,\mathscr E)\) 中的这个元素。

\medskip\noindent\textbf{c) 归结到 \(S\) 为剩余域代数闭的 artin 局部概形的情形。}

由于 \(K_0\) 与 \(J\otimes_{S_0}\mathcal O_{T_0}\) 为凝聚层，\(\mathscr E\) 也是凝聚层；此外，\(\mathscr E\) 的支集包含在 \(T_0\) 中，因为 \(J\otimes_{S_0}\mathcal O_{T_0}\) 具有这一性质。为证明 \(\mathscr E\) 的截面 \(\Psi\) 为零，只须证明它在 \(T_0\) 的每一点 \(u\)，它在 \(\mathscr E\) 在 \(u\) 处的纤维中的像为零。凝聚层的 \(\mathcal Ext^i\) 的构造与平坦的基扩张可交换，\footnote{编者注：参见 EGA \(0_{\mathrm{III}}\)，12.3.5。}故可归结为证明
\[
  T_0\cap\operatorname{Spec}\mathcal O_u
\]
存在一个在 \(S\) 上平坦的提升。设 \(s\) 为 \(u\) 在 \(S\) 中的像。可用 \(\operatorname{Spec}\mathcal O_s\) 替换 \(S\)，用 \(G\times_S\operatorname{Spec}\mathcal O_s\) 替换 \(G\)。

必要时再作一次忠实平坦扩张，可假设 \(\mathcal O_s\) 的剩余域代数闭（EGA \(0_{\mathrm{III}}\)，10.3.1）。设 \(\mathfrak m\) 为 \(\mathcal O_s\) 的极大理想，置
\[
  S_n=\operatorname{Spec}(\mathcal O_s/\mathfrak m^n).
\]
假设对每个 \(n>0\)，已证明将
\[
  T_0\times_S S_n=(T_0)_n
\]
提升为 \(G_n=G\times_S S_n\) 中在 \(S_n\) 上平坦的子概形之障碍消失，设 \(T_n\) 为 \((T_0)_n\) 在 \(S_n\) 上唯一的、同时为 \(G_n\) 的子环面的平坦提升。若 \(n>n'\)，则显然
\[
  T_n\times_{S_n}S_{n'}=T_{n'}.
\]
若 \(u\) 为 \(T_0\) 中位于 \(s\) 之上的一点，将下述引理应用于 \((T_0)_n\cap\operatorname{Spec}\mathcal O_u\) 的一族相容的提升
\[
  T_n\cap\operatorname{Spec}\mathcal O_u,
\]
即可知 \(T_0\cap\operatorname{Spec}\mathcal O_u\) 确有在 \(\mathcal O_s\) 上平坦的提升。因此只须在 \(S=S_n\) 为剩余域代数闭的 artin 局部环的谱时证明 \(\Psi=0\)，并证明下述引理。

\SGAStatement{\SGAInlineRefTarget{sga3:xv:lemma:2-8}引理 2.8}{设 \(A\to B\) 为 noether 局部环之间的局部同态，\(\mathfrak m\) 为 \(A\) 的极大理想，\(J\) 为 \(A\) 的平方为零的理想，\(M\) 为有限型 \(B\)-模，置
\[
  A_0=A/J,\qquad B_0=B/JB,\qquad M_0=M/JM.
\]
设 \(N_0\) 为 \(M_0\) 的一个在 \(A_0\) 上平坦的商 \(B_0\)-模。对每个 \(n>0\)，用 \(A_n,A_{0,n}\) 等表示由基扩张
\[
  A\longrightarrow A/\mathfrak m^n=A_n
\]
得到的对象，用 \(\overline{J}_n\) 表示 \(J\) 在 \(A_n\) 中的像。对每个 \(n>0\)，设 \(N_n\) 为 \(M_n\) 的一个商 \(B_n\)-模，在 \(A_n\) 上平坦并提升 \(N_{0,n}\)，且假设当 \(n\ge n'\) 时，\(N_{n'}\) 由 \(N_n\) 经基扩张 \(A_n\to A_{n'}\) 得到。则存在 \(M\) 的商 \(B\)-模 \(N\)，使其在 \(A\) 上平坦且提升 \(N_0\)。}
\SGARefTarget{sga3:xv:lemma:2:8}

\SGAProof{\SGARefLink{sga3:xv:lemma:2-8}{2.8} 的证明}
设 \(P_0\) 为满态射 \(M_0\to N_0\) 的核。对每个 \(n>0\)，有横行正合的交换图：
\[
\tag{$\ast_n$}\label{sga3:xv:lemma:2-8:equation:ast-n}
\begin{tikzcd}[column sep=2em,row sep=large]
&&&& 0 \arrow[d] & \\
&&&& P_{0,n} \arrow[d] & \\
& 0 \arrow[r]
  & \overline{J}_nM_n \arrow[r] \arrow[dd]
  & M_n \arrow[r] \arrow[dd]
  & M_{0,n} \arrow[r] \arrow[dd] & 0 \\
\overline{J}_n\otimes_{A_{0,n}}M_{0,n}
      \arrow[urr] \arrow[drr] &&&&& \\
& 0 \arrow[r]
  & \overline{J}_n\otimes_{A_{0,n}}N_{0,n} \arrow[r]
  & N_n \arrow[r] \arrow[d]
  & N_{0,n} \arrow[r] \arrow[d] & 0 \\
&&& 0 & 0 &
\end{tikzcd}
\]
此外，由假设，图 \((\ast)_{n+1}\) 模 \(\mathfrak m^n\) 约化后即为 \((\ast)_n\)。

Artin--Rees 引理（Bourbaki，\emph{Algèbre commutative}，第三章，第 3 节，推论 1）表明，分别由态射
\[
\begin{gathered}
  JM\longrightarrow \overline{J}_nM_n,\\
  J\otimes_{A_0}M_0
    \longrightarrow \overline{J}_n\otimes_{A_{0,n}}M_{0,n},\\
  J\otimes_{A_0}N_0
    \longrightarrow \overline{J}_n\otimes_{A_{0,n}}N_{0,n},
\end{gathered}
\]
的核所定义的 \(JM\)、\(J\otimes_{A_0}M_0\)、\(J\otimes_{A_0}N_0\) 上的滤过均为 \(\mathfrak mB\)-良滤过。用 \(\widehat Q\) 表示 \(B\)-模 \(Q\) 关于 \(\mathfrak mB\)-进拓扑的分离完备化。对图 \((\ast)_n\) 取逆极限，得到下列交换图，其中两个横行仍正合：
\[
\tag{$\widehat\ast$}\label{sga3:xv:lemma:2-8:equation:widehat-ast}
\begin{tikzcd}[column sep=normal,row sep=large]
&&&& 0 \arrow[d] & \\
&&&& \widehat P_0 \arrow[d] & \\
& 0 \arrow[r]
  & \widehat{JM} \arrow[r] \arrow[dd]
  & \widehat M \arrow[r] \arrow[dd]
  & \widehat M_0 \arrow[r] \arrow[dd] & 0 \\
\widehat{J\otimes_{A_0}M_0}
      \arrow[urr] \arrow[drr] &&&&& \\
& 0 \arrow[r]
  & \widehat{J\otimes_{A_0}N_0} \arrow[r]
  & \varprojlim_n N_n \arrow[r] \arrow[d]
  & \widehat N_0 \arrow[r] \arrow[d] & 0 \\
&&& 0 & 0 &
\end{tikzcd}
\]
现 \((B,\mathfrak mB)\) 为 Zariski 环，且 \(J\otimes_{A_0}N_0\) 在 \(B\) 上有限型，因而关于 \(\mathfrak mB\)-进拓扑分离。于是图 \((\widehat\ast)\) 表明，由满态射 \(M_0\to N_0\) 诱导的态射
\[
  J\otimes_{A_0}M_0\longrightarrow J\otimes_{A_0}N_0
\]
通过 \(JM\) 分解：
\[
\begin{tikzcd}[column sep=large,row sep=large]
J\otimes_{A_0}M_0 \arrow[rr,"\mathrm{can.}"] \arrow[dr]
  && JM \arrow[dl] \\
& J\otimes_{A_0}N_0 &
\end{tikzcd}
\]

在这些条件下，第 \SGARefLink{sga3:iii:lemma:4-1}{III~4.1} 讲和 \SGARefLink{sga3:iii:proposition:4-3}{4.3} 表明：存在 \(M\) 的一个在 \(A\) 上平坦且提升 \(N_0\) 的商 \(B\)-模 \(N\)，当且仅当某个元素
\[
  g\in E=\operatorname{Ext}^{1}_B
  (P_0,J\otimes_{A_0}N_0)
\]
为零。更确切地，正合列
\[
  0\longrightarrow JM\longrightarrow M
  \longrightarrow M_0\longrightarrow0
\]
定义元素 \(f\in F\)，其中
\[
  F=\operatorname{Ext}^{1}_B(M_0,JM),
\]
且 \(g\) 为 \(f\) 在由
\[
  P_0\longrightarrow M_0,
  \qquad
  JM\longrightarrow J\otimes_{A_0}N_0
\]
通过双函子性引出的自然态射 \(F\to E\) 下的像。图 \((\widehat\ast)\) 与第 \SGARefLink{sga3:iii:lemma:4-1}{III~4.1} 讲表明，\(g\) 在
\[
  \widehat E\simeq
  \operatorname{Ext}^{1}_{\widehat B}
  (\widehat P_0,\widehat{J\otimes_{A_0}N_0})
\]
中的像为零。但 \(E\) 在 \(B\) 上有限型，因而关于 \(\mathfrak mB\)-进拓扑分离。所以 \(g=0\)，\SGARefLink{sga3:xv:lemma:2-8}{2.8} 得证。

% 底本源行 849--1084。
\medskip\noindent\textbf{d) \(\boldsymbol{\mathscr E}\) 的研究。}

现在假设 \(S\) 为 artin 局部环 \(A\) 的谱，且其剩余域 \(k\) 代数闭。设 \(A_0\) 为 \(S_0\) 的环，\(\mathfrak m_0\) 为 \(A_0\) 的极大理想。由于 \(S_0\) 为 artin 概形，\(T_0\) 在 \(G\) 中闭（第 VI\textsubscript{B}~\SGARefLink{cor:sga3-VIB-1.4.2}{1.4.2} 讲）。因此可取 \(U=G\)，于是
\[
  \mathscr E=
  \mathcal Ext^{1}_{\mathcal O_G}
  (K_0,J\otimes_{S_0}\mathcal O_{T_0}).
\]
设 \(B_0\) 为仿射 \(S_0\)-概形 \(T_0=\mathbb G_{m,S_0}\) 的环，\(\overline B_0\) 为其特殊纤维
\[
  T_0\times_{S_0}\operatorname{Spec}(k)=\mathbb G_{m,k}
\]
的环。层 \(\mathscr E\) 为凝聚 \(\mathcal O_{T_0}\)-模，因而由一个有限型 \(B_0\)-模定义，记为 \(E\)。

考虑 \(E\) 关于 \(\mathfrak m_0B_0\)-进滤过的相伴分次模：
\[
  E_r=\mathfrak m_0^rE/\mathfrak m_0^{r+1}E.
\]
每个 \(E_r\) 都是有限型 \(\overline B_0\)-模；由于 \(S_0\) 为 artin 概形，充分大的 \(r\) 都满足 \(E_r=0\)。

设 \(g\) 为 \(G\) 在 \(S\) 上的截面，其在 \(S_0\) 上的限制是 \(T_0\) 的截面 \(g_0\)。就障碍问题而言，由 \(g\) 定义的左平移给出“整个情形的一个自同构”。特别地，\(g\) 诱导层 \(\mathscr E\) 的一个保持障碍 \(\Psi\) 不变的 \(S_0\)-自同构。更确切地，\(g\) 定义 \(B_0\)-模 \(E\) 的一个半线性自同构，其所对应的 \(B_0\) 的 \(A_0\)-自同构由 \(T_0\) 中的 \(g_0\) 平移定义。因此模 \(\mathfrak m_0^{r+1}\) 约化后，\(g\) 定义 \(E_r\) 的一个半线性自同构，其所对应的 \(\overline B_0\) 的 \(k\)-自同构由 \(T_0\times_{S_0}\operatorname{Spec}(k)\) 中的
\[
  g_0\times_{S_0}\operatorname{Spec}(k)
\]
平移定义。

\SGAStatement{引理 2.9}{对每个整数 \(r\ge0\)，\(E_r\) 为局部自由 \(\overline B_0\)-模。}
\SGARefTarget{sga3:xv:lemma:2:9}

设 \(x\) 为 \(T_0\) 的一点，\(\kappa(x)\) 为其剩余域，置
\[
  (E_r)_x=E_r\otimes_{\overline B_0}\kappa(x),
\]
即 \(E_r\) 在 \(x\) 处的“纤维”。用 \(\ell(x)\) 表示 \((E_r)_x\) 在 \(\kappa(x)\) 上的秩，\(\ell\) 表示当 \(x\) 遍取 \(T_0\) 各点时 \(\ell(x)\) 的最大值。设 \(L\) 为
\[
  \operatorname{Spec}\overline B_0=\mathbb G_{m,k}
\]
中使 \(E_r\) 在其上局部自由、秩为 \(\ell\) 的最大闭子概形（TDTE IV，\SGARefLink{sga3:xv:lemma:3:6}{引理 3.6}）。取 \(\beta\in L(k)\)；这样的点存在，因为 \(L\) 在代数闭域 \(k\) 上有限型。再取一点 \(\alpha\in\mathbb G_{m,k}(k)\)，其阶为 \(q\) 的某个幂 \(n\)。点 \(\alpha\) 为 \(k\)-有理点。根据假设，\({}_nT_0\) 可提升为一个在 \(S\) 上平展的子概形 \(M_n\)，故存在 \(G\) 在 \(S\) 上的截面 \(a\)，提升 \(\alpha\)，且其在 \(S_0\) 上的限制为 \(T_0\) 的截面。

前述观察表明，\(E_r\) 在 \(\mathbb G_{m,k}(k)\) 的两点 \(\beta\)、\(\alpha\beta\) 处的纤维 \(k\)-同构。阶为 \(q\) 的幂的点在 \(\mathbb G_{m,k}\) 中稠密，这些点由 \(\beta\) 平移所得的点也稠密。由于 \(L\) 在 \(\mathbb G_{m,k}\) 中闭且 \(\mathbb G_{m,k}\) 约化，有 \(L=\mathbb G_{m,k}\)，故 \(E_r\) 为秩为 \(\ell\) 的局部自由 \(\overline B_0\)-模。

\medskip\noindent\textbf{e) \SGARefLink{sga3:xv:theorem:2:1}{定理 2.1} 证明的完成。}

设 \(K_0(n)\) 为 \(\mathcal O_{G_0}\) 中定义闭子概形 \({}_nT_0\) 的理想层。因此层 \(K_0\)\footnote{编者注：此层在\SGARefLink{sga3:xv:lemma:2-7}{引理 2.7} 之后引入。}为 \(K_0(n)\) 的子层。为简便起见，置
\[
  R=J\otimes_{S_0}\mathcal O_{T_0},
  \qquad
  R(n)=J\otimes_{S_0}\mathcal O_{{}_nT_0}.
\]
层 \(R(n)\) 为 \(R\) 的商，且有图
\[
\begin{tikzcd}[column sep=large,row sep=large]
&& K_0 \arrow[d] \\
&& K_0(n) \arrow[d] \\
0 \arrow[r]
  & J\mathcal O_G \arrow[r] \arrow[d]
  & \mathcal O_G \arrow[r]
  & \mathcal O_{G_0} \arrow[r] & 0 \\
  & R \arrow[d] \\
  & R(n)
\end{tikzcd}
\]
利用 \(\operatorname{Ext}^{1}_{\mathcal O_G}\) 的双函子性，得到交换图
\[
\begin{tikzcd}[column sep=normal,row sep=large]
\mathcal{E}xt^{1}_{\mathcal O_G}
  (\mathcal O_{G_0},J\mathcal O_G)
  \arrow[dr] \\
&
\mathcal{E}xt^{1}_{\mathcal O_G}(K_0(n),R)
  \arrow[r] \arrow[d]
&
\mathcal{E}xt^{1}_{\mathcal O_G}(K_0,R)=\mathscr E
  \arrow[d] \\
&
\mathcal{E}xt^{1}_{\mathcal O_G}(K_0(n),R(n))
  \arrow[r]
&
\mathcal{E}xt^{1}_{\mathcal O_G}(K_0,R(n)).
\end{tikzcd}
\]
仍用 \(\Phi\) 表示正合列
\[
  0\longrightarrow J\mathcal O_G
  \longrightarrow\mathcal O_G
  \longrightarrow\mathcal O_{G_0}
  \longrightarrow0
\]
所定义的 \(\operatorname{Ext}^{1}_{\mathcal O_G}
(\mathcal O_{G_0},J\mathcal O_G)\) 中的元素，于是 \(\Psi\) 为 \(\Phi\) 在 \(\mathscr E\) 中的像。由于 \({}_nT_0\) 可提升为 \(G\) 中在 \(S\) 上平坦的子概形 \(M_n\)，\(\Phi\) 在
\[
  \operatorname{Ext}^{1}_{\mathcal O_G}(K_0(n),R(n))
\]
中的像为零（第 \SGARefLink{sga3:iii:lemma:4-1}{III~4.1} 讲）。故 \(\Phi\) 在 \(\operatorname{Ext}^{1}_{\mathcal O_G}(K_0,R(n))\) 中的像也为零，而这个像同时就是 \(\Psi\) 的像。

\SGAStatement{引理 2.10}{对每个整数 \(n\)，典范态射
\[
  \operatorname{Ext}^{1}_{\mathcal O_G}(K_0,R)
    \otimes_{B_0}\mathcal O_{{}_nT_0}
  \longrightarrow
  \operatorname{Ext}^{1}_{\mathcal O_G}(K_0,R(n))
\]
为单射。}
\SGARefTarget{sga3:xv:lemma:2:10}

事实上，仿射概形 \(T_0=\mathbb G_{m,S_0}\) 的环为
\[
  B_0=A_0[T,T^{-1}].
\]
子概形 \({}_nT_0\) 由
\[
  h(n)=T^n-1
\]
的消失定义。这个截面正则（EGA \(0_{\mathrm{IV}}\)，15.2.2），且在每个基变换 \(S'_0\to S_0\) 之后仍然正则。故有层的正合列
\[
  0\longrightarrow\mathcal O_{T_0}
  \xrightarrow{\,h(n)\,}\mathcal O_{T_0}
  \longrightarrow\mathcal O_{{}_nT_0}
  \longrightarrow0.
\]
由于 \({}_nT_0\) 在 \(S\) 上平坦，在 \(\mathcal O_{S_0}\) 上与 \(J\) 张量得到正合列
\[
  0\longrightarrow R\xrightarrow{\,h(n)\,}R
  \longrightarrow R(n)\longrightarrow0,
\]
从而得到 Ext 群的正合列
\[
  \cdots\longrightarrow
  \operatorname{Ext}^{1}_{\mathcal O_G}(K_0,R)
  \xrightarrow{\,h(n)\,}
  \operatorname{Ext}^{1}_{\mathcal O_G}(K_0,R)
  \longrightarrow
  \operatorname{Ext}^{1}_{\mathcal O_G}(K_0,R(n)),
\]
引理得证。

以上表明，对每个等于 \(q\) 的幂的整数 \(n\)，\(\Psi\) 在 \(\mathscr E/h(n)\mathscr E\) 中的像为零。为证明 \(\Psi=0\)，只须证明
\[
  \Psi\in\mathfrak m_0^r\mathscr E
  \quad\Longrightarrow\quad
  \Psi\in\mathfrak m_0^{r+1}\mathscr E.
\]
设 \(\overline\Psi\) 为 \(\Psi\) 在 \(E_r\) 中的像。存在元素 \(\Psi(n)\in\mathscr E\) 使
\[
  \Psi=\Psi(n)h(n)
  \qquad(n\text{ 为 }q\text{ 的幂}).
\]
前已注意到，\(h(n)\) 在 \(\overline B_0\) 中的像 \(\overline{h(n)}\) 仍为 \(\overline B_0\)-正则元。由 \SGARefLink{sga3:xv:lemma:2:9}{2.9}，对每个 \(s\)，\(E_s\) 都在 \(\overline B_0\) 上局部自由，故 \(E_s\) 上乘以 \(\overline{h(n)}\) 是单射。立即得到 \(\Psi(n)\in\mathfrak m_0^r\mathscr E\)。用 \(\overline{\Psi(n)}\) 表示它在 \(E_r\) 中的像；于是
\[
  \overline\Psi=\overline{h(n)}\,\overline{\Psi(n)}
  \qquad(n\text{ 为 }q\text{ 的幂}).
\]
因此，\(\mathbb G_{m,k}\) 中 \(\overline\Psi\) 消失的点所成集合 \(F\)，包含阶为 \(q\) 的幂的点所成的稠密集。此外，\(F\) 闭，因为 \(E_r\) 在 \(\overline B_0\) 上局部自由，且 \(\mathbb G_{m,k}\) 约化。因此 \(\overline\Psi=0\)。

\SGARefLink{sga3:xv:theorem:2:1}{定理 2.1} 的证明至此完成。

% 底本源行 1085--1338。
\SGASection{3}{由底集刻画子环面}
\SGARefTarget{sga3:xv:section:3}

\SGASubsection{3.1}{定理的陈述}
\SGARefTarget{sga3:xv:subsection:3:1}

\noindent\textbf{记号。}---若 \(X\) 为预概形，用 \(\operatorname{ens}(X)\) 表示 \(X\) 的底集。若 \(X\)、\(S'\) 为两个 \(S\)-预概形，
\[
  X_{S'}=X\times_S S',
  \qquad u:X_{S'}\longrightarrow X
\]
为典范态射，且 \(E\) 为 \(\operatorname{ens}(X)\) 的子集，则用 \(E_{S'}\)（或 \(E\times_S S'\)）表示 \(\operatorname{ens}(X_{S'})\) 的子集 \(u^{-1}(E)\)。

\SGAStatement{\SGAInlineRefTarget{sga3:xv:theorem:3-1}定理 3.1}{设 \(S\) 为局部 noether 预概形，\(G\) 为有限型 \(S\)-群预概形，\(q>0\) 为在 \(S\) 上可逆的整数，\(E\) 为 \(\operatorname{ens}(G)\) 的子集。考虑关于 \(E\) 的下列断言：
\begin{enumerate}
\item[\normalfont\SGAInlineRefTarget{sga3:xv:theorem:3-1:item:i}(i)] 集合 \(E\) 为 \(G\) 的一个子环面 \(T\) 的底集。
\item[\normalfont\SGAInlineRefTarget{sga3:xv:theorem:3-1:item:ii}(ii)]
  \begin{enumerate}
  \item[\normalfont a)] 对 \(S\) 的每一点 \(s\)，存在 \(G_s\) 的子环面 \(T_s\)，使
  \[
    E\cap\operatorname{ens}(G_s)=E_s
  \]
  为 \(T_s\) 的底集。
  \item[\normalfont b)] 当 \(n\) 遍取 \(q\) 的幂时，存在 \(G\) 的一族相容的（参见 \SGARefLink{sga3:xv:corollary:2-2}{2.2}）乘法型子群概形 \(M_n\)，使对 \(S\) 的每一点 \(s\)，
  \[
    (M_n)_s={}_nT_s.
  \]
  \end{enumerate}
\item[\normalfont\SGAInlineRefTarget{sga3:xv:theorem:3-1:item:iii}(iii)]
  \begin{enumerate}
  \item[\normalfont a)] 同 \textup{\SGARefLink{sga3:xv:theorem:3-1:item:ii}{(ii)} a)}。
  \item[\normalfont b)] 集合 \(E\) 在 \(\operatorname{ens}(G)\) 中局部闭，且 \(E\) 在 \(S\) 上的纤维的维数局部常值。
  \end{enumerate}
\item[\normalfont\SGAInlineRefTarget{sga3:xv:theorem:3-1:item:iv}(iv)]
  \begin{enumerate}
  \item[\normalfont a)] 同 \textup{\SGARefLink{sga3:xv:theorem:3-1:item:ii}{(ii)} a)}。
  \item[\normalfont b)] 对每个为剩余域代数闭的完备离散赋值环之谱的 \(S\)-概形 \(S'\)，\(E_{S'}\) 为 \(G_{S'}\) 的一个子环面的底集。
  \end{enumerate}
\end{enumerate}
则有下列蕴涵关系：
\begin{enumerate}
\item[\normalfont A)] \(\textup{(i)}\Longleftrightarrow
\textup{(ii)}\Longrightarrow\textup{(iii)}
\Longrightarrow\textup{(iv)}\)。
\item[\normalfont B)] 若 \(G\) 在 \(S\) 上分离，则 \(\textup{(iii)}\Longleftrightarrow\textup{(iv)}\)，且 \(E\) 为闭集。
\item[\normalfont C)] 在下列两种情形，条件 \textup{\SGARefLink{sga3:xv:theorem:3-1:item:i}{(i)}}、\textup{\SGARefLink{sga3:xv:theorem:3-1:item:ii}{(ii)}}、\textup{\SGARefLink{sga3:xv:theorem:3-1:item:iii}{(iii)}} 等价（若 \(G\) 在 \(S\) 上分离，则也与 \textup{\SGARefLink{sga3:xv:theorem:3-1:item:iv}{(iv)}} 等价）：
  \begin{enumerate}
  \item[\normalfont 第一种情形：]
    \begin{enumerate}
    \item[\normalfont a)] 预概形 \(S\) 约化，或 \(G\) 在 \(S\) 上平坦；且
    \item[\normalfont b)] 对 \(S\) 的每一点 \(s\)，\(T_s\) 为 \(G_s\) 的中心环面。
    \end{enumerate}
  \item[\normalfont 第二种情形：] \(S\) 正规。
  \end{enumerate}
\end{enumerate}
此外，在上述两种情形，底集为 \(E\) 的环面 \(T\) 唯一。}
\SGARefTarget{sga3:xv:theorem:3:1}

\SGAStatement{\SGAInlineRefTarget{sga3:xv:remark:3-2}注 3.2}{\SGAInlineRefTarget{sga3:xv:remark:3-2:item:a}\textup{a)} 当 \(S\) 约化时，在 \textup{\SGARefLink{sga3:xv:theorem:3-1:item:ii}{(ii)}} 中无须假设族 \(M_n\) 相容，因为这一条件可由其他假设推出。事实上，若 \(m\) 整除 \(n\)，子群 \({}_mM_n\)、\(M_m\) 均在 \(S\) 上平展，从而约化；它们具有相同的底空间，因而相同。

\textup{b)} 若不再假设 \(S\) 正规，则 \(\textup{(iii)}\Rightarrow\textup{(i)}\) 一般不再成立，即使 \(S\) 约化、几何单分支且 \(G\) 为 \(S\) 上的光滑群概形也是如此。事实上，考虑 \(\operatorname{SL}_{2,S}\) 中由形如
\[
  \begin{pmatrix}t&u\\0&t^{-1}\end{pmatrix}
\]
的矩阵构成的 Borel 子群，其中 \(S\) 为域 \(k\) 上环为
\[
  k[x,y]/(y^2-x^3)
\]
的仿射曲线。如下定义 \(E\)。在 \(S\) 的尖点，即 \(x=y=0\) 处，取对角环面 \(u=0\)。在补开集 \(x\ne0\) 上，取对角环面经
\[
  \begin{pmatrix}1&y/x\\0&1\end{pmatrix}
\]
共轭得到的环面。所得集合 \(E\) 满足 \textup{\SGARefLink{sga3:xv:theorem:3-1:item:iii}{(iii)} a)}；它在 \(G\) 中还为闭集，以 \(E\) 为底集的约化子概形由方程
\[
  \begin{cases}
    xu+y(t-t^{-1})=0,\\
    u^2-x(t-t^{-1})^2=0
  \end{cases}
\]
定义。因此它不是 \(G\) 的子环面，因为其在尖点上的纤维不约化。}
\SGARefTarget{sga3:xv:remarks:3:2}

\noindent\textbf{\SGARefLink{sga3:xv:theorem:3-1}{3.1} 的证明安排。}第一部分将证明下列“容易”的蕴涵关系：
\[
\begin{tikzcd}[column sep=huge,row sep=large]
 & \textup{(ii)} \arrow[dr,Rightarrow] & \\
\textup{(i)} \arrow[ur,Rightarrow] \arrow[dr,Rightarrow] && \textup{(iv)} \\
 & \textup{(iii)} \arrow[ur,Rightarrow] &
\end{tikzcd}
\]
以及
\[
  \textup{(iv)}
  \Longrightarrow
  [\,\textup{(iii)}\text{ 且 }E\text{ 闭}\,]
  \quad\text{若 \(G\) 在 \(S\) 上分离}.
\]
其余较为细致的蕴涵关系将分三步证明：
\begin{enumerate}
\item[\normalfont I)] 在 C) 的第一种情形的假设下，将 \(\textup{(iii)}\Rightarrow\textup{(i)}\) 先归结到 \(S\) 正规的情形。
\item[\normalfont II)] 当 \(S\) 正规时，\(\textup{(iii)}\Rightarrow\textup{(ii)}\)。
\item[\normalfont III)] \(\textup{(ii)}\Rightarrow\textup{(i)}\)。
\end{enumerate}

\SGASubsection{3.2}{\SGARefLink{sga3:xv:theorem:3-1}{3.1} 中“容易”断言的证明}
\SGARefTarget{sga3:xv:subsection:3:2}

蕴涵 \(\textup{(i)}\Rightarrow\textup{(ii)}\)、\(\textup{(i)}\Rightarrow\textup{(iii)}\) 显然。

\medskip
\noindent\(\textup{(iii)}\Rightarrow\textup{(iv)}\)。
用 \(S'\) 替换 \(S\) 后，可假设 \(S\) 为离散赋值环的谱。设 \(t\) 为 \(S\) 的一般点，\(s\) 为其闭点。

由于 \(E\) 局部闭，存在 \(G\) 中一个以 \(E\) 为底空间的约化子预概形，仍记为 \(E\)。由 \textup{\SGARefLink{sga3:xv:theorem:3-1:item:iii}{(iii)} a)}，其一般纤维 \(E_t\) 即 \(G_t\) 的子环面 \(T_t\)，设其秩为 \(r\)。用 \(E'\) 表示 \(E_t\) 在 \(E\) 中的概形闭包。预概形 \(E'\) 不可约，其闭纤维 \(E'_s\) 非空（包含 \(G_s\) 的单位截面）；因此 \(E'_s\) 的维数为 \(r\)（EGA IV~14.3.10）。但 \(E'_s\) 是 \(E_s\) 的闭子概形，而由 \textup{\SGARefLink{sga3:xv:theorem:3-1:item:iii}{(iii)} b)}，\(E_s\) 的维数为 \(r\)，且其底空间与 \(T_s\) 相同，故不可约。因此 \(E'_s\)、\(E_s\) 有相同的底空间，进而 \(E'=E\)。这证明 \(E\) 在 \(S\) 上平坦。

现在设 \(E''\) 为 \(E_t\) 在 \(G\) 中的概形闭包。则 \(E''\) 为 \(G\) 中一个在 \(S\) 上平坦的子群预概形（第 \SGARefLink{sga3:viii:lemma:7:1}{VIII~7.1} 讲），它包含 \(E'\)，因而包含 \(E\)。其闭纤维 \(E''_s\) 为 \(G_s\) 的一个维数为 \(r\) 的代数子群（同上）。由于 \(T_s\) 为 \(G_s\) 的一个维数为 \(r\) 的闭不可约子集，\(E_s\) 与 \(E''_s\) 的单位连通分支 \((E''_s)^0\) 具有相同的底集。用 \((E'')^0\) 表示 \(E''\) 的“单位连通分支”，即 \(E''\) 的一个开子群，其补集为 \(E''_s\) 中不含原点的不可约分支之并。于是
\[
  E=\operatorname{ens}\bigl((E'')^0\bigr).
\]
由于 \(E\)、\(E''\) 约化，甚至有 \(E=(E'')^0\)。最后，\(E\) 为 \(G\) 中一个在 \(S\) 上平坦、有限型且纤维连通的子群预概形，因而分离（第 VI\textsubscript{B}~\SGARefLink{sga3:xv:theorem:5-2}{5.2} 讲）；其一般纤维为环面 \(T_t\)，其约化闭纤维为环面 \(T_s\)。故 \(E\) 是环面（第 \SGARefLink{sga3:x:theorem:8-8}{X~8.8} 讲）。

\medskip
\noindent\(\textup{(ii)}\Rightarrow\textup{(iv)}\)。
同样可假设 \(S\) 为离散赋值环的谱，并沿用上述记号。\(T_t\) 在 \(G\) 中的概形闭包 \(T\) 为 \(G\) 中一个在 \(S\) 上平坦的子群预概形，且对每个为 \(q\) 的幂的 \(n\)，它都包含 \(M_n\)。因此，其闭纤维是 \(G_s\) 中包含所有 \((M_n)_s\) 的闭子概形，从而包含 \(T_s\)。根据维数，\(T\) 的“单位连通分支”\(T^0\) 的底集为 \(E\)，且前述论证表明 \(T^0\) 为 \(G\) 的子环面。

\medskip
\noindent\(\textup{(iv)}\Rightarrow
[\,\textup{(iii)}\text{ 且 }E\text{ 闭}\,]\)，若 \(G\) 在 \(S\) 上分离。下面证明 \(E\) 闭。\footnote{编者注：在证明过程中将得到蕴涵 \(\textup{(iv)}\Rightarrow\textup{(iii)}\)。}先证明下述引理。

\SGAStatement{引理 3.3}{若 \SGARefLink{sga3:xv:theorem:3-1}{3.1} \textup{\SGARefLink{sga3:xv:theorem:3-1:item:iv}{(iv)}} 中的条件成立，则 \(E\) 为 \(\operatorname{ens}(G)\) 的局部可构造子集。}
\SGARefTarget{sga3:xv:lemma:3:3}

按通常方法，可归结到 \(S\) 为 noether 整预概形的情形，设其一般点为 \(\eta\)。缩小 \(S\) 后，可假设（第 VI\textsubscript{B}~\SGARefLink{par:sga3-VIB-10.10}{10.10} 讲）存在 \(G\) 的子群概形 \(H\)，在 \(S\) 上平坦且纤维连通，并使其一般纤维 \(H_\eta\) 为 \(T_\eta\)。为证明\SGARefLink{sga3:xv:lemma:3:3}{引理 3.3}，只须证明 \(\operatorname{ens}(H)=E\)。

设 \(s\) 为 \(S\) 的一点。由 EGA II~7.1.9，存在 \(S\)-概形 \(S'\)，它为离散赋值环的谱，且可假设该环完备、剩余域代数闭，其一般点 \(t'\) 映至 \(\eta\)，闭点 \(s'\) 映至 \(s\)。由 \textup{\SGARefLink{sga3:xv:theorem:3-1:item:iv}{(iv)} b)}，存在 \(G_{S'}\) 中以 \(E_{S'}\) 为底空间的子环面 \(T_{S'}\)。\(G_{S'}\) 的两个子群预概形 \(T_{S'}\)、\(H_{S'}=H\times_S S'\) 均在 \(S'\) 上平坦、纤维连通且一般纤维同为 \(T_{t'}\)。因此两者都等于 \(T_{t'}\) 在 \(G_{S'}\) 中的概形闭包的单位连通分支。故 \(\operatorname{ens}(H_s)=E_s\)，从而 \(\operatorname{ens}(H)=E\)，引理得证。

已知 \(E\) 局部可构造，为证明其闭，只须证明 \(G\) 中 \(E\) 的点的每个特殊化仍属于 \(E\)。用通常的方法，可归结到 \(S\) 为剩余域代数闭的完备离散赋值环之谱的情形。此时，由 \textup{\SGARefLink{sga3:xv:theorem:3-1:item:iv}{(iv)} b)} 给出的底空间为 \(E\) 的 \(G\) 的子环面在 \(G\) 中闭，因为 \(G\) 分离（第 \SGARefLink{sga3:viii:proposition:7:12}{VIII~7.12} 讲）。

% 底本源行 1339--1721。
\SGASubsection{3.3}{\SGARefLink{sga3:xv:theorem:3-1}{3.1} 证明的继续}
\SGARefTarget{sga3:xv:subsection:3:3}

\noindent\textbf{I) 在 C 的第一种情形，将 \(\textup{(iii)}\Rightarrow\textup{(i)}\) 归结到 \(S\) 正规的情形。}

\smallskip
\noindent\textit{a) 归结到 \(S\) 仿射且约化的情形。}\quad
我们假设对 \(S\) 的每一点 \(s\)，\(E_s\) 为 \(G_s\) 的一个中心子环面的底空间。\(T\) 的唯一性由第 \SGARefLink{sga3:ix:theorem:5:1:bis}{IX~5.1 bis} 讲推出，且 \(T\) 必然包含在中心中（同上）。由此唯一性，为证明 \(T\) 的存在性，可以假设 \(S\) 为 Noether 仿射概形，设其环为 \(A\)。若 \(S\) 不约化，则根据假设，\(G\) 在 \(S\) 上平坦。由 \SGARefLink{sga3:xv:corollary:2:3}{2.3}，此时只须对 \(S_{\mathrm{red}}\) 与 \(G\times_S S_{\mathrm{red}}\) 解决问题。因此还可假设 \(S\) 约化。

\smallskip
\noindent\textit{b) 归结到 \(A\) 为 \(\mathbb Z\) 上有限型代数的情形。}\quad
我们将用到如下引理。

\SGAStatement{引理 3.4}{设 \(k\) 为域，\(G\) 为代数 \(k\)-群，\(E\) 为 \(G\) 的子集。设 \(k'/k\) 为域扩张，并假设 \(E_{k'}\) 为 \(G_{k'}\) 的一个子环面 \(T_{k'}\) 的底集。若 \(T_{k'}\) 包含在 \(G_{k'}\) 的中心中，或 \(k\) 为完美域，则 \(E\) 为 \(G\) 的一个子环面 \(T\) 的底集。}
\SGARefTarget{sga3:xv:lemma:3:4}

\SGAProof{证明}
由 fpqc 下降，只须证明 \(T_{k'}\) 在 \(G_{k''}\) 中的两个逆像相同，其中
\[
  k''=k'\otimes_k k'.
\]
它们有相同的底集。若 \(T_{k'}\) 包含在中心中，则由第 \SGARefLink{sga3:ix:theorem:5:1:bis}{IX~5.1 bis} 讲，两者相同。若 \(k\) 完美，则 \(k''\) 约化；这两个逆像在 \(k''\) 上光滑，因而约化，所以相同。

\SGAStatement{\SGAInlineRefTarget{sga3:xv:remark:3-5}注 3.5}{\textup{(1)} 特别地，当 \(S\) 的剩余域均完美时，可应用\SGARefLink{sga3:xv:lemma:3:4}{引理 3.4}。

\textup{\SGARefLink{sga3:xv:theorem:5:2}{(2)}} 在\SGARefLink{sga3:xv:theorem:3-1}{定理 3.1} C 的第一种情形的假设下，对每个如 \SGARefLink{sga3:xv:theorem:3-1}{3.1} \textup{\SGARefLink{sga3:xv:theorem:3-1:item:iv}{(iv)} b)} 中所述的 \(S\)-概形 \(S'\)，\(G_{S'}\) 中底空间为 \(E_{S'}\) 的子环面包含在中心中。}
\SGARefTarget{sga3:xv:remarks:3:5}

\SGAStatement{引理 3.6}{设 \(S\) 为仿射概形的一个滤过逆系统 \(S_i\) 的投射极限，\(H\) 为有限型乘法型 \(S\)-群。则存在指标 \(i\)、\(S_i\) 上的乘法型群 \(H_i\) 以及 \(S\)-同构
\[
  H_i\times_{S_i}S\xrightarrow{\sim}H.
\]
若 \(H\) 为等平凡的，则可选取等平凡的 \(H_i\)。}
\SGARefTarget{sga3:xv:lemma:3:6}

\SGAProof{证明}
由于 \(H\) 在 \(S\) 上有限表现，存在指标 \(i\)、在 \(S_i\) 上有限表现的 \(S_i\)-群预概形 \(H_i\) 以及 \(S\)-同构
\[
  H_i\times_{S_i}S\xrightarrow{\sim}H
\]
（第 VI\textsubscript{B}~\SGARefLink{prop:sga3-VIB-10.16}{10.16} 讲）。增大 \(i\) 后，可假设 \(H_i\) 为乘法型（第 \SGARefLink{sga3:ix:proposition:2:11}{IX~2.11} 讲）。尚须证明可选取拟等平凡的 \(H_i\)，并在 \(H\) 等平凡时使其等平凡。

因此可假设 \(H\) 被一个有限表现的平展满态射 \(S'\to S\) 平凡化。当 \(i\) 充分大时，存在 \(S_i\)-预概形 \(S'_i\)、有限表现的平展满态射 \(S'_i\to S_i\) 以及 \(S\)-同构
\[
  S'_i\times_{S_i}S\xrightarrow{\sim}S'
\]
（EGA IV~17.16）。对充分大的 \(j\)，置
\[
  S'_j=S'_i\times_{S_i}S_j,\qquad
  H'_j=H_j\times_{S_j}S'_j,\qquad
  H'=H\times_S S'.
\]
由 \(S'\) 的选取方式，存在有限生成 Abel 群 \(M\) 以及群概形的一个 \(S'\)-同构
\[
  D_{S'}(M)\xrightarrow{\sim}H'.
\]
由于各 \(S'_i\) 拟紧且 \(S'=\varprojlim S'_i\)，第 VI\textsubscript{B}~10 讲表明，存在指标 \(j\) 以及群概形的一个 \(S'_j\)-同构
\[
  D_{S'_j}(M)\xrightarrow{\sim}H'_j.
\]
这恰好说明 \(H_j\) 为拟等平凡的乘法型群。当 \(H\) 等平凡时，以有限平展态射 \(S'\to S\) 对 \(H\) 的平凡化作同样论证即可。

现在可完成归约 b)。\(S\) 的环 \(A\) 为其各个在 \(\mathbb Z\) 上有限型的子环 \(A_i\) 的滤过直极限。置 \(S_i=\operatorname{Spec}A_i\)，设
\[
  u_j:S\longrightarrow S_j,\qquad
  u_{jk}:S_k\longrightarrow S_j\quad(k\geq j)
\]
为过渡态射。由于 \(S\) 为 Noether 概形且 \(G\) 在 \(S\) 上有限表现，存在指标 \(i\) 以及在 \(S_i\) 上有限型的 \(S_i\)-群预概形 \(G_i\)，使 \(G\) 与 \(G_i\times_{S_i}S\) \(S\)-同构。同样，由于 \(E\) 在 \(G\) 中局部闭，可以假设存在 \(G_i\) 的局部闭子集 \(E_i\) 使
\[
  E=E_i\times_{S_i}S
\]
（EGA IV~8.3.11）。对 \(j>i\)，置
\[
  G_j=G_i\times_{S_i}S_j,\qquad
  E_j=E_i\times_{S_i}S_j,
\]
并用 \(Q_j\) 表示 \(S_j\) 中满足下述条件的点 \(s\) 所成的集合：\((E_j)_s\) 为 \((G_j)_s\) 的一个中心子环面的底集。

由 \SGARefLink{sga3:xv:lemma:3:4}{3.4} 可得
\[
  Q_k=u_{jk}^{-1}(Q_j)\quad(k\geq j),
\]
且根据假设，
\[
  \operatorname{ens}(S)=u_j^{-1}(Q_j)\quad(j\geq i).
\]
我们进一步断言，\(Q_j\) 是归纳可构造（ind-constructible）的（EGA IV~1.9.4）。事实上，由于 \(S_j\) 为 Noether 概形，只须证明下述事实（EGA IV~1.9.10）：设 \(S\) 为一般点为 \(\eta\) 的 Noether 整概形，且 \(E_\eta\) 为 \(G_\eta\) 的一个中心子环面 \(T_\eta\) 的底集，则存在 \(\eta\) 的邻域 \(U\)，使每个 \(s\in U\) 的 \(E_s\) 都有相同性质。缩小 \(S\) 后，可以假设 \(G\) 的中心 \(Z\) 可表，且存在 \(G\) 的子群概形 \(T\)，其一般纤维为 \(T_\eta\)（第 VI\textsubscript{B} 讲，第 10 节）。\(\operatorname{ens}(G)\) 的两个局部闭子集 \(\operatorname{ens}(T)\)、\(E\) 在一般纤维上相同。因此可找到 \(\eta\) 的邻域 \(U\)，使在 \(U\) 上有
\[
  \operatorname{ens}(T)=E
\]
（EGA IV~8.3.11）。出于同样理由，可以假设 \(T\) 在 \(U\) 上包含在中心中；由 3.6，它是环面。

既然 \(Q_j\) 归纳可构造，EGA IV~8.3.4 表明，对充分大的 \(j\)，有 \(Q_j=\operatorname{ens}(S_j)\)。此时以 \(S_j,G_j,E_j\) 替换 \(S,G,E\)，即归结到 \(S\) 为 \(\mathbb Z\) 上有限型仿射约化概形的情形。

\smallskip
\noindent\textit{c) 归结到 \(S\) 为完备约化 Noether 局部环的谱的情形。}\quad
由于底空间为 \(E\) 的环面唯一，用通常的方法（EGA IV~8）及引理 3.6，可将 \(S\) 替换为其某点处局部环 \(A\) 的谱。设 \(\widehat S\) 为 \(A\) 关于其极大理想定义的拓扑的完备化 \(\widehat A\) 的谱。由于 \(A\) 是 \(\mathbb Z\) 上有限型代数的局部化，\(\widehat S\) 仍约化（EGA IV~7.6.5）。

我们断言，只须在基扩张 \(\widehat S\to S\) 后解决问题。事实上，设 \(\widehat T\) 为 \(G_{\widehat S}\) 中底空间为 \(E_{\widehat S}\) 的子环面。它在 \(G_{\widehat S\times_S\widehat S}\) 中的两个逆像均为中心子环面，且底空间相同，因而相同（第 \SGARefLink{sga3:ix:theorem:5:1:bis}{IX~5.1 bis} 讲）。由 fpqc 下降，\(\widehat T\) 来自 \(G\) 的一个环面 \(T\)，它就解决了问题（参见 \SGARefLink{sga3:xv:lemma:3:4}{3.4}）。

\smallskip
\noindent\textit{d) 一个下降引理。}\quad
回顾 TDTE I 中指出的有限态射的下述性质。设 \(S,S'\) 为两个预概形，\(u:S'\to S\) 为有限态射。则：
\begin{enumerate}
\item[\normalfont\SGAInlineRefTarget{sga3:xv:lemma:3:6:item:1}(1)] 态射 \(u\) 为满态射，当且仅当典范层态射
\[
  \mathcal O_S\longrightarrow u_*(\mathcal O_{S'})
\]
为单射。
\item[\normalfont(2)] 态射 \(u\) 为有效满态射（第 \SGARefLink{sga3:iv:definition:1:3}{IV~1.3} 讲），当且仅当典范图
\[
  0\longrightarrow\mathcal O_S
  \longrightarrow u_*(\mathcal O_{S'})
  \rightrightarrows
  u_*(\mathcal O_{S'})\otimes_{\mathcal O_S}u_*(\mathcal O_{S'})
\]
正合。
\item[\normalfont\SGAInlineRefTarget{sga3:xv:lemma:3:6:item:3}(3)] 若 \(S\) 还是 Noether 预概形，且 \(u\) 为满态射，则 \(u\) 为有限个有限有效满态射的复合。
\end{enumerate}

现在可证明如下引理，其证明用到了非平坦下降的方法。

\SGAStatement{\SGAInlineRefTarget{sga3:xv:lemma:3-7}引理 3.7}{设 \(S\) 为局部 Noether 预概形，\(G\) 为有限型 \(S\)-群预概形，\(S'\) 为预概形，\(u:S'\to S\) 为有限态射。对每个 \(S\)-预概形 \(T\)，用 \(\mathrm M(T)\) 表示 \(G_T\) 的乘法型子群的集合\textup{（分别地，\(G_T\) 的中心乘法型子群的集合）}，则 \(\mathrm M\) 自然地成为 \(\operatorname{Sch}/S\) 上的一个反变函子。若 \(u\) 为有效满态射\textup{（分别地，满态射）}，则集合图
\[
  \tag{*}\label{sga3:xv:lemma:3-7:equation:star}
  \mathrm M(S)\longrightarrow\mathrm M(S')
  \rightrightarrows\mathrm M(S'\times_S S')
\]
正合。}
\SGARefTarget{sga3:xv:lemma:3:7}

\SGAProof{证明}
\textit{i) 归结到 \(u\) 为有效满态射的情形。}\quad
这里考虑 \(G\) 的中心乘法型子群函子。\(\mathrm M(S)\to\mathrm M(S')\) 的单射性在 \(S\) 上是局部性质；一旦有此单射性，\((*)\) 的正合性在 \(S\) 上也是局部性质。因此可假设 \(S\) 为 Noether 仿射概形。

考虑 \(u:S'\to S\) 为两个有限满态射的复合的情形：
\[
  S'\xrightarrow{v}S''\xrightarrow{w}S.
\]
我们断言，若以下两个图
\[
  \tag{*'}\label{sga3:xv:lemma:3-7:equation:star-prime}
  \mathrm M(S'')\longrightarrow\mathrm M(S')
  \rightrightarrows\mathrm M(S'\times_{S''}S')
\]
及
\[
  \tag{*''}\label{sga3:xv:lemma:3-7:equation:star-prime-prime}
  \mathrm M(S)\longrightarrow\mathrm M(S'')
  \rightrightarrows\mathrm M(S''\times_S S'')
\]
正合，则 \((*)\) 也正合。

事实上，\(\mathrm M(S)\to\mathrm M(S')\) 的单射性显然。若 \(T'\) 属于
\[
  \operatorname{Ker}\bigl(
    \mathrm M(S')\rightrightarrows
    \mathrm M(S'\times_S S')\bigr),
\]
则当然也属于
\[
  \operatorname{Ker}\bigl(
    \mathrm M(S')\rightrightarrows
    \mathrm M(S'\times_{S''}S')\bigr).
\]
由 \((*')\) 的正合性，\(T'\) 来自 \(\mathrm M(S'')\) 的唯一元素 \(T''\)。只须证明 \(T''\) 属于
\[
  \operatorname{Ker}\bigl(
    \mathrm M(S'')\rightrightarrows
    \mathrm M(S''\times_S S'')\bigr),
\]
因为此时 \((*'')\) 的正合性使 \(T''\) 下降到 \(S\)。设 \(T''_1,T''_2\) 为 \(T''\) 在 \(\mathrm M(S''\times_S S'')\) 中的两个逆像。它们都是 \(G_{S''\times_S S''}\) 的中心乘法型子群，故为证明 \(T''_1=T''_2\)，只须验证它们有相同的纤维（第 \SGARefLink{sga3:ix:theorem:5:1:bis}{IX~5.1 bis} 讲）。考虑交换图
\[
\begin{tikzcd}[column sep=huge,row sep=large]
S' \arrow[d,"v"'] &
  S'\times_S S' \arrow[l,shift left=.55ex]
    \arrow[l,shift right=.55ex] \arrow[d,"v\times v"]\\
S'' &
  S''\times_S S'' . \arrow[l,shift left=.55ex]
    \arrow[l,shift right=.55ex]
\end{tikzcd}
\]
态射 \(v\) 为有限满态射，因而由上述性质 \textup{\SGARefLink{sga3:xv:lemma:3:6:item:1}{(1)}} 知其占优，又为闭态射，所以满射；因此 \(v\times v\) 也满射。设 \(x''\) 为 \(S''\times_S S''\) 的一点，\(x'\) 为 \(S'\times_S S'\) 中位于 \(x''\) 之上的一点。图的交换性表明，\((T''_1)_{x''}\)、\((T''_2)_{x''}\) 在 \(G_{x'}\) 中的两个逆像相同。因此
\[
  (T''_1)_{x''}=(T''_2)_{x''}
\]
（EGA IV~2.2.15）。

用上述论证，并对由性质 \textup{\SGARefLink{sga3:xv:lemma:3:6:item:3}{(3)}} 给出的将 \(u:S'\to S\) 分解为有效满态射之复合的因子个数作直接归纳，即知为证明中心乘法型子群的 \((*)\) 正合性，只须处理 \(u\) 为有效满态射的情形。最后，再应用一次第 \SGARefLink{sga3:ix:theorem:5:1:bis}{IX~5.1 bis} 讲可知，只须在 \(\mathrm M\) 为 \(G\) 的所有乘法型子群所成函子且 \(u\) 为有效满态射时证明 \SGARefLink{sga3:xv:lemma:3-7}{3.7}。

\smallskip
\noindent\textit{ii) \(\mathrm M(S)\to\mathrm M(S')\) 的单射性。}\quad
由于 \(u:S'\to S\) 为满态射，典范态射
\[
  \mathcal O_S\longrightarrow u_*(\mathcal O_{S'})
\]
为单射。此外，乘法型 \(S\)-群在 \(S\) 上平坦。因此 \(\mathrm M(S)\to\mathrm M(S')\) 的单射性由下述更一般的引理得出。

\SGAStatement{\SGAInlineRefTarget{sga3:xv:lemma:3-8}引理 3.8}{设 \(f:X\to S\)、\(g:S'\to S\) 为预概形的态射，\(\mathcal F\) 为拟凝聚 \(\mathcal O_X\)-模，置
\[
  X'=X\times_S S',
  \qquad
  \mathcal F'=\mathcal F\otimes_{\mathcal O_S}\mathcal O_{S'}.
\]
用 \(\mathrm Q(\mathcal F)\) 表示 \(\mathcal F\) 的拟凝聚且在 \(S\) 上平坦的商 \(\mathcal O_X\)-模 \(\mathcal G\) 的集合；相对于 \(\mathcal F',X'\) 和 \(S'\)，类似地定义 \(\mathrm Q(\mathcal F')\)。假设 \(g\) 拟紧且
\[
  \mathcal O_S\longrightarrow g_*(\mathcal O_{S'})
\]
为单射。则典范映射
\[
  \begin{aligned}
  \mathrm Q(\mathcal F)&\longrightarrow\mathrm Q(\mathcal F'),\\
  \mathcal G&\longmapsto
    \mathcal G\otimes_{\mathcal O_S}\mathcal O_{S'}
  \end{aligned}
\]
为单射。}
\SGARefTarget{sga3:xv:lemma:3:8}

\SGAProof{证明}
事实上，可假设 \(S\) 仿射，再假设 \(X\) 仿射。由于 \(S'\) 拟紧，它有有限仿射开覆盖。用这些开子集的不交和替换 \(S'\)，即可假设 \(S'\) 仿射。于是归结到下述引理，其证明是直接的。

\SGAStatement{引理 3.9}{设 \(A\to A'\) 为单的环同态，\(M\) 为 \(A\)-模，\(N=M/P\) 为在 \(A\) 上平坦的商 \(A\)-模，置
\[
  M'=M\otimes_A A',
  \qquad
  N'=N\otimes_A A'=M'/P'.
\]
则 \(P\) 为 \(P'\) 在典范同态 \(M\to M'\) 下的逆像。因此 \(N\)、\(P\) 由 \(N'\)、\(P'\) 决定。}
\SGARefTarget{sga3:xv:lemma:3:9}

% 底本源行 1722--2099。
\smallskip
\noindent\textit{iii) \SGARefLink{sga3:xv:lemma:3-7}{3.7} \((*)\) 在 \(\mathrm M(S')\) 处的正合性。}\quad
设
\[
  T'\in\operatorname{Ker}\bigl(
    \mathrm M(S')\rightrightarrows
    \mathrm M(S'\times_S S')\bigr).
\]
假设当 \(S\) 为 Noether 局部环的谱时，iii) 已得证。则对 \(S\) 的每一点 \(s\)，存在 \(G\times_S\operatorname{Spec}(\mathcal O_{S,s})\) 的乘法型子群 \(T_s\)，由
\[
  T'\times_{S'}u^{-1}\!\bigl(\operatorname{Spec}(\mathcal O_{S,s})\bigr)
\]
下降而来。由第 VI\textsubscript{B} 讲第 10 节和 \SGARefLink{sga3:xv:lemma:3:6}{3.6}，存在 \(s\) 在 \(S\) 中的开邻域 \(U\)，以及 \(G\) 在 \(U\) 上的乘法型子群概形 \(T\)，延拓 \(T_s\)。设 \(U'\) 为 \(U\) 在 \(S'\) 中的逆像。此时 \(G'|_{U'}\) 的两个子概形 \(T'|_{U'}\)、\(T\times_U U'\) 在 \(u^{-1}(\operatorname{Spec}(\mathcal O_{S,s}))\) 上相同。

将后者视为 \(u^{-1}(V)\) 的投射极限，其中 \(V\) 遍取 \(s\) 在 \(U\) 中的开邻域，EGA IV~8 表明，存在包含 \(s\) 的开子集 \(V\subset U\)，使 \(T\times_U U'\)、\(T'\) 在其逆像上相同。故在所述假设下，\(T'\) 可在 \(S\) 上局部下降；由 ii) 已证的唯一性，它因而可整体下降。总之，只须在 \(S\) 为 Noether 局部环的谱时证明 iii)。

设 \(\widehat S\) 为 \(S\) 的环的完备化之谱，置
\[
  S''=S'\times_S S',
  \qquad
  \widehat S'=\widehat S\times_S S',
  \qquad
  \widehat S''=\widehat S\times_S S''
      =\widehat S'\times_S S'.
\]
我们断言，只须证明
\[
  \tag{\(\widehat *\)}\label{sga3:xv:lemma:3-9:equation:widehat-star}
  \mathrm M(\widehat S)\longrightarrow
  \mathrm M(\widehat S')
  \rightrightarrows
  \mathrm M(\widehat S'')
\]
在 \(\mathrm M(\widehat S')\) 处正合。这由下列交换图得出：
\[
\begin{tikzcd}[column sep=large,row sep=large]
\mathrm M(S) \arrow[r] \arrow[d]
  & \mathrm M(S') \arrow[r,shift left=.55ex]
      \arrow[r,shift right=.55ex] \arrow[d]
  & \mathrm M(S'') \arrow[d]\\
\mathrm M(\widehat S) \arrow[r]
      \arrow[d,shift left=.55ex] \arrow[d,shift right=.55ex]
  & \mathrm M(\widehat S') \arrow[r,shift left=.55ex]
      \arrow[r,shift right=.55ex]
      \arrow[d,shift left=.55ex] \arrow[d,shift right=.55ex]
  & \mathrm M(\widehat S'')\\
\mathrm M(\widehat S\times_S\widehat S)
  \arrow[r,"f"]
  & \mathrm M(\widehat S'\times_{S'}\widehat S').
\end{tikzcd}
\]
其中，第二行根据假设在 \(\mathrm M(\widehat S')\) 处正合，前两列由 fpqc 下降正合；将 ii) 应用于从 \(S'\to S\) 经平坦基变换 \(\widehat S\times_S\widehat S\to S\) 得到的有限满态射
\[
  \widehat S'\times_{S'}\widehat S'
  \longrightarrow
  \widehat S\times_S\widehat S,
\]
即得 \(f\) 为单射。图追踪留给读者。d)\SGARefLink{sga3:xv:theorem:5:2}{(2)} 中回顾的有限有效满态射之刻画表明，从 \(S'\to S\) 经平坦基变换 \(\widehat S\to S\) 得到的 \(\widehat S'\to\widehat S\) 仍为有限有效满态射。因此还可归结到 \(S\) 为完备 Noether 局部环的谱的情形。

设 \(S_0\) 为 \(S\) 中以闭点为底空间的约化闭子概形。设 \(T'\) 如上，用 \(T''\) 表示其在 \(\mathrm M(S'')\) 中的公共像。由于
\[
  S'_0=S'\times_S S_0
\]
在 \(S_0\) 上忠实平坦，由 fpqc 下降，存在 \(G_{S_0}\) 的乘法型子群 \(T_0\)，其在 \(G_{S'_0}\) 中的逆像为 \(T'_{S'_0}\)。由于 \(S\) 为完备 Noether 局部概形，存在乘法型 \(S\)-群 \(T\) 及 \(S_0\)-态射
\[
  u_0:T\times_S S_0\longrightarrow T_0
\]
（第 \SGARefLink{sga3:x:corollary:3:3}{X~3.3} 讲）。仍由第 \SGARefLink{sga3:x:corollary:3:3}{X~3.3} 讲，\(u_0\) 在 \(S'_0\) 上的逆像唯一延拓为 \(S'\)-同构
\[
  u':T_{S'}\xrightarrow{\sim}T'.
\]
此处应注意，由于 \(S'\) 在完备局部概形 \(S\) 上有限，它为有限个完备局部概形的不交和。

\(u'\) 在 \(S''\) 上的两个逆像均为从 \(T_{S''}\) 到 \(T''\) 的态射，它们在 \(S_0\times_S S''\) 上相同，因而由同一结果，在处处相同。由于 \(T\) 在 \(S\) 上平坦，且 \(S'\to S\) 为有限有效满态射，TDTE I 第 8 页给出正合图
\[
  \operatorname{Hom}_S(T,G)\longrightarrow
  \operatorname{Hom}_{S'}(T_{S'},G_{S'})
  \rightrightarrows
  \operatorname{Hom}_{S''}(T_{S''},G_{S''}).
\]
因此 \(u'\) 来自预概形的一个态射 \(u:T\to G\)。同样，\(T\times_S T\) 在 \(S\) 上平坦，故
\[
  \operatorname{Hom}_S(T\times_S T,G)\longrightarrow
  \operatorname{Hom}_{S'}(T_{S'}\times_{S'}T_{S'},G_{S'})
\]
为单射。由于 \(u'\) 为群同态，立即可知 \(u\) 也是群同态。

此外，\(u\) 为单态射。先注意 \(\operatorname{Ker}(u)\) 在 \(S\) 上平坦。为证明这一点，可以假设 \(S\) 为 Artin 局部概形（EGA \(0_{\mathrm{III}}\)，10.2.2），从而 \(G\) 分离（第 VI\textsubscript{A} 讲，\SGARefLink{sga3:VIA:0.3}{0.3}）；于是 \(\operatorname{Ker}(u)\) 为乘法型（第 \SGARefLink{sga3:ix:theorem:6:8}{IX~6.8} 讲），因而在 \(S\) 上平坦。由于
\[
  \operatorname{Ker}(u)\times_S S'
    =\operatorname{Ker}(u')=0,
\]
\SGARefLink{sga3:xv:lemma:3-8}{引理 3.8} 给出 \(\operatorname{Ker}(u)=0\)。最后，单态射 \(u\) 为浸入（\SGARefLink{sga3:viii:remark:7:13}{第 VIII 讲，注 7.13}），其像 \(u(T)\) 为 \(\mathrm M(S)\) 的一个元素，在 \(\mathrm M(S')\) 中的像为 \(T'\)。\SGARefLink{sga3:xv:lemma:3-7}{引理 3.7} 得证。

\smallskip
\noindent\textit{e) I) 的证明的结束。}\quad
由归约 c)，可归结到 \(S\) 为完备约化 Noether 局部环 \(A\) 的谱的情形。设 \(S'\) 为 \(A\) 的正规化 \(A'\) 的谱。Nagata 定理表明，\(A'\) 在 \(A\) 上有限（EGA \(0_{\mathrm{IV}}\)，23.1.5），且由于 \(A\to A'\) 为单射，\(S'\to S\) 为满态射。假设 \(G_{S'}\) 有环面 \(T'\)，其底空间为 \(E_{S'}\)。\(T'\) 在 \(G_{S'\times_S S'}\) 中的两个逆像为底空间相同的中心子环面，因而相同（第 IX~5 bis 讲）。故由\SGARefLink{sga3:xv:lemma:3-7}{引理 3.7}，\(T'\) 来自 \(G_S\) 的一个中心子环面 \(T\)，显然其底空间为 \(E\)。因此只须证明 \(T'\) 的存在性。这就归结到正规情形，并完成了 I) 的证明。

\medskip
\noindent\textit{II) 当 \(S\) 正规时，证明 \textup{\SGARefLink{sga3:xv:theorem:3-1:item:iii}{(iii)}} 蕴涵 \textup{\SGARefLink{sga3:xv:theorem:3-1:item:ii}{(ii)}}。}\quad
可以只考虑 \(S\) 整的情形，设其一般点为 \(t\)。对每个为 \(q\) 的幂的整数 \(n\)，用 \({}_nG\) 表示 \(G\) 中取 \(n\) 次幂的“核”。由 \SGARefLink{sga3:xv:theorem:3-1:item:iii}{(iii)}(b)，\(E\) 在 \(G\) 中局部闭，故 \(\operatorname{ens}({}_nG)\) 与 \(E\) 的交在 \(\operatorname{ens}(G)\) 中局部闭。用 \(E(n)\) 表示 \(G\) 中底空间为 \(\operatorname{ens}({}_nG)\cap E\) 的约化子概形。

下面证明结构态射
\[
  E(n)\longrightarrow S
\]
分离且泛开。这两种性质均有赋值判据（EGA II~7.2.3 和 EGA IV~14.5.8）。因此设 \(S'\) 为 \(S\)-概形，它是剩余域代数闭的完备离散赋值环的谱。已证 \SGARefLink{sga3:xv:theorem:3-1:item:iii}{3.1(iii)} 蕴涵 \SGARefLink{sga3:xv:theorem:3-1:item:iv}{3.1(iv)}，故存在 \(G_{S'}\) 中底空间为 \(E_{S'}\) 的子环面 \(T'\)。而 \({}_nT'\) 在 \(S'\) 上有限平展，因而在 \(S'\) 上分离且泛开；所以底空间相同的 \(E(n)\times_S S'\) 也有这两种性质。

此外，\(E(n)\) 的所有纤维均有相同数目的几何点，即若 \(T_t\) 的秩为 \(r\)，则几何点数为 \(n^r\)。\footnote{法文版本印为 \(r^n\)；秩为 \(r\) 的环面的 \(n\)-挠点数应为 \(n^r\)。}最后，一般纤维 \(E(n)_t\) 约化，因而等于 \({}_nT_t\)，从而在 \(\operatorname{Spec}(t)\) 上平展。由于 \(S\) 正规，由 SGA I~10.11 可知 \(E(n)\) 是 \(S\) 的平展覆盖。

若 \(s\) 为 \(S\) 的一点，则 \(E(n)_s\) 在 \(\operatorname{Spec}\kappa(s)\) 上平展，从而约化，故等于乘法型群 \({}_nT_s\)。下面证明 \(E(n)\) 为 \(G\) 的子群概形。设
\[
  m:E(n)\times_S E(n)\longrightarrow G
\]
为 \(G\) 的乘法诱导的态射。底映射 \(\operatorname{ens}(m)\) 通过 \(E(n)\) 分解；由于源 \(E(n)\times_S E(n)\) 在 \(S\) 上平展，从而约化，\(m\) 本身也通过概形 \(E(n)\) 分解。由第 \SGARefLink{sga3:x:corollary:4-8:item:a}{X~4.8(a)} 讲，\(E(n)\) 是乘法型子群。如 \SGARefLink{sga3:xv:remark:3-2:item:a}{3.2(a)} 已指出，子群族 \(E(n)\) 必然相容。这证明了当 \(S\) 正规时，\SGARefLink{sga3:xv:theorem:3-1:item:iii}{(iii)} 蕴涵 \SGARefLink{sga3:xv:theorem:3-1:item:ii}{(ii)}。

\medskip
\noindent\textit{III) 证明 \textup{\SGARefLink{sga3:xv:theorem:3-1:item:ii}{(ii)}} 蕴涵 \textup{\SGARefLink{sga3:xv:theorem:3-1:item:i}{(i)}}。}\quad
事实上，我们将证明存在唯一的 \(G\) 的子环面 \(T\)，其底空间为 \(E\)，且对每个为 \(q\) 的幂的 \(n\)，
\[
  {}_nT=M_n
\]
成立。唯一性立即由第 \SGARefLink{sga3:ix:corollary:4:8:item:b}{IX~4.8(b)} 讲得出。为证明存在性，利用此唯一性，可以依次假设：
\begin{enumerate}
\item[\normalfont a)] \(S\) 为 Noether 预概形。
\item[\normalfont b)] 各 \(M_n\) 为中心子群。事实上，根据 \SGARefLink{sga3:xv:lemma:2:5}{2.5} 和 \SGARefLink{sga3:xv:lemma:2-5-bis}{2.5 bis}，只须对充分大的 \(n\)，以
\[
  Z_n=\operatorname{Centr}_G(M_n)
\]
替换 \(G\)。
\item[\normalfont c)] \(S\) 为局部环的谱。假设在每个基变换 \(\operatorname{Spec}(\mathcal O_{S,s})\to S\)（\(s\in S\)）之后问题都已解决，用 \(T_s\) 表示所得到的 \(G\times_S\operatorname{Spec}(\mathcal O_{S,s})\) 的子环面。由第 VI\textsubscript{B} 讲第 10 节和 \SGARefLink{sga3:xv:lemma:3:6}{3.6}，存在 \(s\) 的开邻域 \(U\)，以及 \(G|_U\) 的子环面 \(T\)，延拓 \(T_s\)。已证 \SGARefLink{sga3:xv:theorem:3-1:item:ii}{3.1(ii)} 蕴涵 \SGARefLink{sga3:xv:theorem:3-1:item:iv}{3.1(iv)}；由于 \(S\) 为 Noether 预概形，故由 \SGARefLink{sga3:xv:lemma:3:3}{3.3}，\(E\) 可构造。因此缩小 \(U\) 后，可假设
\[
  \operatorname{ens}(T)=E\times_S U
\]
（EGA IV~9.5.2）。对每个为 \(q\) 的幂的 \(n\)，群 \({}_nT\)、\(M_n\times_S U\) 均为 \(G|_U\) 的乘法型子群，且有相同的纤维；由于 \(M_n\) 为中心子群，由第 \SGARefLink{sga3:ix:corollary:5:3:bis}{IX~5.3 bis} 讲，两者相同。因此问题在 \(S\) 上有局部解，唯一性使这些局部解组成整体解。
\item[\normalfont d)] \(S\) 为完备 Noether 局部环的谱，且若 \(s\) 为闭点，则 \(T_s\) 分裂。这由 EGA \(0_{\mathrm{III}}\)，10.3.1 及 fpqc 下降得出。
\item[\normalfont e)] \(S\) 约化。应用 \SGARefLink{sga3:xv:corollary:2-2}{2.2} 即可。
\item[\normalfont f)] \(S\) 正规。应用 Nagata 有限性定理（EGA \(0_{\mathrm{IV}}\)，23.1.5）以及对中心子环面的引理 3.7 即可。
\end{enumerate}

经过这些归约，\(T_s\) 分裂，从而 \((M_n)_s\) 分裂。由第 \SGARefLink{sga3:x:corollary:3:3}{X~3.3} 讲，\(M_n\) 本身分裂。若 \(t\) 为 \(S\) 的任一点，则 \(T_t\) 的所有子群 \({}_nT_t\)（\(n\) 为 \(q\) 的幂）都分裂；于是第 \SGARefLink{sga3:x:proposition:1:4}{X~1.4} 讲立即表明 \(T_t\) 本身分裂。只须再证明下述引理。

\SGAStatement{引理 3.10}{在 \SGARefLink{sga3:xv:theorem:3-1:item:ii}{3.1(ii)} 的假设下，进一步假设 \(S\) 为 Noether 正规预概形，且对 \(S\) 的每一点 \(s\)，\(T_s\) 为分裂环面。则存在唯一的 \(G\) 的子环面 \(T\)，其底空间为 \(E\)，且对每个为 \(q\) 的幂的 \(n\)，
\[
  {}_nT=M_n
\]
成立。此外，\(T\) 分裂。}
\SGARefTarget{sga3:xv:lemma:3:10}

\SGAProof{证明}
由于 \(T\) 在 \(S\) 上光滑，因而约化，故唯一性成立。为证明存在性，可假设 \(S\) 不可约，设其一般点为 \(\eta\)。设 \(T_\eta\) 的秩为 \(r\)，置
\[
  T'=\mathbb G_{m,S}^{\,r},
\]
并选取同构
\[
  u_\eta:T'_\eta\xrightarrow{\sim}T_\eta.
\]
用 \({}_nu_\eta\) 表示其在 \({}_nT'_\eta\) 上的限制。由于 \({}_nT'\) 与 \(M_n\) 分裂，\({}_nu_\eta\) 唯一延拓为 \(S\)-同构
\[
  {}_nu:{}_nT'\xrightarrow{\sim}M_n.
\]
我们断言，对 \(S\) 的每一点 \(s\)，存在必然唯一的群同构
\[
  u_s:T'_s\xrightarrow{\sim}T_s
\]
延拓每个 \(({}_nu)_s\)，其中 \(n\) 为 \(q\) 的幂。

事实上，设 \(S^1\) 为 \(S\)-概形，它为剩余域代数闭的完备离散赋值环的谱，其一般点 \(t_1\) 映至 \(\eta\)，闭点 \(s_1\) 映至 \(s\)（EGA II~7.1.9）。由 3.1(ii) 蕴涵 3.1(iv) 的证明，存在 \(G_{S^1}\) 的子环面 \(T^1\)，使对每个为 \(q\) 的幂的 \(n\)，
\[
  (M_n)_{S^1}={}_nT^1
\]
成立。由于 \(S^1\) 正规，由第 \SGARefLink{sga3:x:theorem:5-16}{X~5.16} 讲，\(T^1\) 等平凡。等平凡环面的分类（第 \SGARefLink{sga3:x:corollary:1:2}{X~1.2} 讲）与 SGA I 第 V~8.2 讲，再结合 \(T^1\) 的一般纤维分裂，表明 \(T^1\) 分裂。因此可以将
\[
  u_\eta\times_\eta t_1
\]
延拓为群概形的一个 \(S^1\)-同构
\[
  u^1:T'\times_S S^1\xrightarrow{\sim}T^1.
\]
在 \({}_nT'\times_S S^1\) 上，态射 \(u^1\)、\({}_nu\times_S S^1\) 在一般纤维上相同，因而处处相同。将 \(u^1\) 限制到闭纤维，即在域扩张 \(\kappa(s)\to\kappa(s_1)\) 之后给出所需的 \(u_s\)；由第 \SGARefLink{sga3:ix:corollary:4:8:item:a}{IX~4.8(a)} 讲及 fpqc 下降，它可下降到 \(\kappa(s)\)。

概形 \(T'\) 在正规概形 \(S\) 上光滑，因而正规。由 EGA IV\textsubscript{4}，20.4.6 中的有理映射延拓判据，可知存在 \(S\)-态射
\[
  u:T'\longrightarrow G
\]
使其在每个纤维 \(T'_s\) 上的限制为复合态射
\[
  T'_s\xrightarrow{u_s}T_s\longrightarrow G_s.
\]
事实上，对每个上述类型的基变换 \(S^1\to S\)，前述构造恰好给出了所需的在 \(S^1\) 上的延拓。

尚须证明 \(u\) 为同态。分别用 \(m_{T'}\)、\(m_G\) 表示 \(T'\)、\(G\) 中的乘法。
\[
  m_G\circ(u\times_Su)
  \quad\text{与}\quad
  u\circ m_{T'}
\]
的重合子概形包含每个纤维，因为每个 \(u_s\) 都是同态。因此其底空间与 \(T'\times_S T'\) 相同，而后者在 \(S\) 上光滑，从而约化，故重合子概形等于 \(T'\times_S T'\)。最后，\(u\) 在每个纤维上都是单态射，故为单态射；因此它是浸入（第 \SGARefLink{sga3:viii:proposition:7:9}{VIII~7.9} 讲）。\(T'\) 在 \(u\) 下的像是 \(G\) 的一个底集为 \(E\) 的子环面。

\SGARefLink{sga3:xv:theorem:3-1}{定理 3.1} 的证明至此完成。

% 底本源行 2100--2257。
\SGASection{4}{由子群 \({}_nT\) 刻画子环面 \(T\)}
\SGARefTarget{sga3:xv:section:4}

\SGASubsection{4.1}{主要定理的陈述}
\SGARefTarget{sga3:xv:subsection:4:1}

\SGAStatement{\SGAInlineRefTarget{sga3:xv:theorem:4-1}定理 4.1}{设 \(S\) 为连通局部 Noether 预概形，\(G\) 为在 \(S\) 上有限型的 \(S\)-群预概形，\(q>1\) 为在 \(S\) 上可逆的整数，\(r\) 为正整数。对每个为 \(q\) 的幂的整数 \(n\)，设 \(M(n)\) 为 \(G\) 的乘法型子群概形，型为 \((\mathbb Z/n\mathbb Z)^r\)。假设：
\begin{enumerate}
\renewcommand{\labelenumi}{\textup{\alph{enumi})}}
\item 族 \(M(n)\) 相容：若 \(m\mid n\)，则
\[
  {}_mM(n)=M(m);
\]
\item\SGAInlineRefTarget{sga3:xv:theorem:4-1:item:b} 存在点 \(s\in S\) 及 \(G_s\) 的子环面 \(T_s\)，使
\[
  M(n)_s={}_nT_s\qquad\text{对每个 }n;
\]
\item 对每一点 \(t\in S\)，存在 \(G_t\) 的闭仿射子概形 \(F_t\)，包含所有的 \(M(n)_t\)。
\end{enumerate}
则存在唯一的 \(G\) 的子环面 \(T\)，使对每个为 \(q\) 的幂的 \(n\)，有 \({}_nT=M(n)\)。}
\SGARefTarget{sga3:xv:theorem:4:1}

对态射的提升也有类似的定理。

\SGAStatement{\SGAInlineRefTarget{sga3:xv:theorem:4-1-bis}定理 4.1 bis}{设 \(S,G,q\) 如上，\(T\) 为 \(S\)-环面，对每个为 \(q\) 的幂的 \(n\)，设 \(u(n):{}_nT\to G\) 为 \(S\)-同态。假设：
\begin{enumerate}
\renewcommand{\labelenumi}{\textup{\alph{enumi})}}
\item 族 \(u(n)\) 相容：若 \(m\mid n\)，则
\[
  u(m)=u(n)|_{{}_mT};
\]
\item 存在点 \(s\in S\) 及同态 \(u_s:T_s\to G_s\)，使
\[
  u_s|_{{}_nT_s}=u(n)_s
  \qquad\text{对每个为 }q\text{ 的幂的 }n;
\]
\item 对每一点 \(t\in S\)，存在 \(G_t\) 的闭仿射子概形 \(F_t\)，包含所有的 \(u(n)_t({}_nT_t)\)。
\end{enumerate}
则存在唯一的同态 \(u:T\to G\)，其在每个 \({}_nT\) 上的限制均等于 \(u(n)\)，其中 \(n\) 为 \(q\) 的幂。}
\SGARefTarget{sga3:xv:theorem:4:1:bis}

\SGAStatement{注 4.2}{利用离散赋值环的谱上有限型平坦群预概形的 Abel 秩的下半连续性（参见第 \SGARefLink{sga3:x:remark:8-7}{X~8.7} 讲），可弱化\SGARefLink{sga3:xv:theorem:4-1}{定理 4.1} 和 \SGARefLink{sga3:xv:theorem:4-1-bis}{4.1 bis} 的条件 (c)，只要求对 \(S\) 的每个极大点 \(t\) 存在闭仿射子概形 \(F_t\)。}
\SGARefTarget{sga3:xv:remark:4:2}

下面由\SGARefLink{sga3:xv:theorem:4-1}{定理 4.1} 推出\SGARefLink{sga3:xv:theorem:4-1-bis}{定理 4.1 bis}。置 \(G'=G\times_S T\)。对每个为 \(q\) 的幂的 \(n\)，考虑同态
\[
  v(n):{}_nT\longrightarrow G'
\]
使其向 \(G\)、\(T\) 的投影分别为 \(u(n)\) 和典范浸入 \({}_nT\to T\)。于是 \(v(n)\) 为浸入，用 \(M(n)\) 表示其像子群。族 \(M(n)\) 在\SGARefLink{sga3:xv:theorem:4-1}{定理 4.1} 的意义下相容。群 \(M(n)_s\) 等于 \({}_nT'_s\)，其中 \(T'_s\) 是 \(G'_s\) 中由 \(u_s\) 的图给出的子环面；且对每个 \(t\in S\)，\(G'_t\) 的闭仿射子概形 \(F_t\times_tT_t\) 包含所有的群 \(M(n)_t\)。因此由\SGARefLink{sga3:xv:theorem:4-1}{定理 4.1}，存在 \(G'\) 的子环面 \(T'\)，使对每个 \(n\)，\({}_nT'=M(n)\)。

设 \(f:T'\to T\) 为投影 \(G'\to T\) 的限制，\(f(n)\) 为其在 \(M(n)\) 上的限制。则
\[
  f(n)\circ v(n)=\operatorname{id}_{{}_nT}.
\]
纤维 \(T'_s\) 为 \(u_s\) 的图（第 \SGARefLink{sga3:ix:corollary:4:8:item:b}{IX~4.8(b)} 讲），故 \(f_s\) 为同构。群 \(\ker f\)、\(\operatorname{coker}f\) 为乘法型（第 \SGARefLink{sga3:ix:proposition:2:7}{IX~2.7} 讲），且由于 \(S\) 连通，它们的型恒定。因此它们为平凡群，\(f\) 为同构。若 \(v=f^{-1}\)，则
\[
  v|_{{}_nT}=v(n).
\]
将 \(v\) 与投影 \(G'\to G\) 复合，即得所需同态 \(u:T\to G\)。这证明了存在性；唯一性由第 \SGARefLink{sga3:ix:corollary:4:8:item:a}{IX~4.8(a)} 讲推出。

\SGASubsection{4.2}{应用}
\SGARefTarget{sga3:xv:subsection:4:2}

下面推广定理 \SGARefLink{sga3:ix:theorem:7:1}{IX~7.1}。设 \(A\) 为完备 Noether 局部环，\(I\) 为其极大理想，置
\[
  S=\operatorname{Spec}A,\qquad
  S_m=\operatorname{Spec}(A/I^m).
\]
对每个预概形 \(X\)，置 \(X_m=X\times_SS_m\)。

设 \(G\) 为有限型 \(S\)-群预概形，\(T\) 为 \(S\)-环面，\(q\) 为在 \(S\) 上可逆的整数，
\[
  u_m:T_m\longrightarrow G_m\qquad(m\in\mathbb N)
\]
为相容的同态族。当 \(n\) 遍取 \(q\) 的幂时，用 \(u_m(n)\) 表示 \(u_m\) 在 \({}_nT_m\) 上的限制，用
\[
  u(n):{}_nT\longrightarrow G
\]
表示对所有 \(m\) 延拓各态射 \(u_m(n)\) 的唯一同态（\SGARefLink{sga3:xv:proposition:1-6:item:a}{命题 1.6(a)}）。若对每一点 \(t\in S\)，存在 \(G_t\) 的闭仿射子预概形 \(F_t\)，使其对每个为 \(q\) 的幂的 \(n\) 均包含 \(u(n)_t({}_nT_t)\)，则称族 \((u_m)_{m\in\mathbb N}\) \emph{可容许}。后面将看到，这一性质与 \(q\) 无关。

\SGAStatement{命题 4.3}{沿用上述记号，典范映射
\[
 \operatorname{Hom}_{S\text{-}\mathrm{gr}}(T,G)
 \longrightarrow
 \varprojlim_m\operatorname{Hom}_{S_m\text{-}\mathrm{gr}}(T_m,G_m)
\]
将其源等同于靶中由可容许相容族构成的子集。}
\SGARefTarget{sga3:xv:proposition:4:3}

事实上，取 \(s\) 为 \(S\) 的闭点，应用\SGARefLink{sga3:xv:theorem:4-1-bis}{定理 4.1 bis} 即可。

\SGAStatement{推论 4.4}{沿用上述记号，进一步假设 \(G\) 的纤维仿射。则典范映射
\[
 \operatorname{Hom}_{S\text{-}\mathrm{gr}}(T,G)
 \longrightarrow
 \varprojlim_m\operatorname{Hom}_{S_m\text{-}\mathrm{gr}}(T_m,G_m)
\]
为同构。}
\SGARefTarget{sga3:xv:corollary:4:4}

事实上，当 \(G\) 的纤维仿射时，每个相容族都可容许。

\SGAStatement{注 4.5}{当 \(G\) 分离时，命题 4.3 中的环 \(A\) 可替换为 Noether \(I\)-进环：用 EGA III~5.4.1 替代命题 1.6(a) 即可。}
\SGARefTarget{sga3:xv:remark:4:5}
